% !TEX program = pdflatex

\begin{filecontents}[overwrite]{dec_refs.bib}
% refs.bib as of the referee's read --- add them there, whereupon (refs.bib loading first)
% these locals become shadowed duplicates and may be removed.
@book{drescher2006good, author={Drescher, Gary L.}, title={Good and Real: Demystifying Paradoxes from Physics to Ethics}, publisher={MIT Press}, year={2006}}
@book{bostrom2002anthropic, author={Bostrom, Nick}, title={Anthropic Bias: Observation Selection Effects in Science and Philosophy}, publisher={Routledge}, year={2002}}
@comment{EDIT D1 (Fable 5, Aug 29 2026): the levinstein2020cheating, yudkowsky2010tdt, and wolpert2011lesson entries were transcribed from the deposited reference list by the referee; eyeball at merge. Formerly visible note fields.}
@article{levinstein2020cheating, author={Levinstein, Benjamin A. and Soares, Nate}, title={Cheating Death in Damascus}, journal={The Journal of Philosophy}, volume={117}, number={5}, pages={237--266}, year={2020}, doi={10.5840/jphil2020117516}}
% referee's read and should be added there at merge.
@misc{macdermott2023characterising, author={MacDermott, Matt and Everitt, Tom and Belardinelli, Francesco}, title={Characterising Decision Theories with Mechanised Causal Graphs}, year={2023}, howpublished={arXiv:2307.10987}}
@misc{AethusSemantics, author={Benander, Ajax}, title={Counterfactual Evaluation by Reopening Information States}, year={2026}, howpublished={Preprint}, url={https://philpapers.org/rec/BENCEB-2}}
@misc{AethusBrief, author={Benander, Ajax}, title={Aethic Reasoning in Brief: From a Counterfactual Semantics to the Quantum Measurement Problem}, year={2026}, howpublished={Preprint}, url={https://philpapers.org/rec/BENARI-3}, eprint={https://philsci-archive.pitt.edu/id/eprint/30439}, eprinttype={Alternative URL}}
@misc{Nexus, howpublished={Manuscript}, author={Benander, Ajax}, title={Nexic Reasoning: Defining a Generalized Calculus Over Anthropic Parameters}, year={2024}, month = {12}, url={https://philpapers.org/rec/BENNRD}, doi={10.31219/osf.io/jvkcp}, note={Manuscript; the anthropic wing's master, successive deposits December 2024--2026}}



% VERIFY preferred citation form at merge (working paper).
@misc{hanson1998greatfilter,
  author = {Hanson, Robin},
  title  = {The Great Filter --- Are We Almost Past It?},
  year   = {1998}
}

% --- neighbor entries added during drafting; ALL VERIFY-flagged (standard works, details from memory) ---
% Verified Oct 2026 (PhilPapers record): Synthese 187(1), 95--122.
@article{spohn2012reversing,
  author={Spohn, Wolfgang}, title = {Reversing 30 Years of Discussion: Why Causal Decision Theorists Should One-Box},
  journal = {Synthese}, volume = {187}, number = {1}, pages = {95--122}, year = {2012}}
% Volume/issue confirmed Oct 2026 (Phil. Review 121(4), 2012); pages 483--538 from recollection, VERIFY at submission.
@article{price2012causation,
  author = {Price, Huw}, title = {Causation, Chance, and the Rational Significance of Supernatural Evidence},
  journal = {Philosophical Review}, volume = {121}, number = {4}, pages = {483--538}, year = {2012}}
% Fields from recollection, VERIFY at submission: Phil. Review 73(3) (1964), 338--359.
@article{dummett1964bringing,
  author = {Dummett, Michael}, title = {Bringing About the Past},
  journal = {The Philosophical Review}, volume = {73}, number = {3}, pages = {338--359}, year = {1964}}
% Fields from recollection, VERIFY at submission: J. Phil. 78(6) (1981), 331--356.
@article{horgan1981counterfactuals,
  author = {Horgan, Terence}, title = {Counterfactuals and Newcomb's Problem},
  journal = {The Journal of Philosophy}, volume = {78}, number = {6}, pages = {331--356}, year = {1981}}
% Drafting-model recollection: Routledge 1996; eyeball once at merge.
@book{leslie1996end,
  author = {Leslie, John}, title = {The End of the World: The Science and Ethics of Human Extinction},
  publisher = {Routledge}, year = {1996}}
% Drafting-model recollection: 363, 315--319; confirm once at merge.
@article{gott1993implications,
  author = {Gott, III, J. Richard}, title = {Implications of the Copernican Principle for Our Future Prospects},
  journal = {Nature}, volume = {363}, pages = {315--319}, year = {1993}}
% VERIFY at merge: fill the deposit URL once the causality extraction is deposited.
@misc{AethusCausality, author={Benander, Ajax}, title={Causality from Indexed Information States: Screening, Intervention, and the {Bell} Obstruction}, year={2026}, howpublished={Preprint}, url={https://philpapers.org/rec/BENCWT-4}}
% VERIFY at merge (fields from memory): The Direction of Time, UC Press 1956.
@book{reichenbach1956direction,
  author = {Reichenbach, Hans},
  title = {The Direction of Time},
  publisher = {University of California Press}, year = {1956}}
% Fields from recollection, VERIFY at submission: Phil. Review 116(1) (2007), 93--114.
@article{egan2007some,
  author = {Egan, Andy},
  title = {Some Counterexamples to Causal Decision Theory},
  journal = {The Philosophical Review}, volume = {116}, number = {1}, pages = {93--114}, year = {2007}}
% VERIFY at merge (fields from memory): the tickle-defense home.
@book{eells1982rational,
  author = {Eells, Ellery},
  title = {Rational Decision and Causality},
  publisher = {Cambridge University Press}, year = {1982}}
% VERIFY at merge (fields from memory): the medical-Newcomb discussions.
@book{ahmed2014evidence,
  author = {Ahmed, Arif},
  title = {Evidence, Decision and Causality},
  publisher = {Cambridge University Press}, address = {Cambridge}, year = {2014}
}
@article{arntzenius2008no,
  author = {Arntzenius, Frank},
  title = {No regrets, or: {Edith} {Piaf} revamps decision theory},
  journal = {Erkenntnis},
  volume = {68}, number = {2}, pages = {277--297}, year = {2008}
}
@article{joyce2012regret,
  author = {Joyce, James M.},
  title = {Regret and instability in causal decision theory},
  journal = {Synthese},
  volume = {187}, number = {1}, pages = {123--145}, year = {2012}
}
@book{skyrms1990dynamics,
  author = {Skyrms, Brian},
  title = {The Dynamics of Rational Deliberation},
  publisher = {Harvard University Press}, address = {Cambridge, MA}, year = {1990}
}
@book{skyrms1980causal,
  author = {Skyrms, Brian},
  title = {Causal Necessity: A Pragmatic Investigation of the Necessity of Laws},
  publisher = {Yale University Press}, year = {1980}}
% VERIFY at merge: fill the deposit URL and confirm the final title once
% the anthropic extraction is deposited; placeholder fields until then.
@misc{NexicFoundation, author={Benander, Ajax}, title={The Copernican Method Against the Copernican Principle: Anthropic Typicality over Partial Information States}, year={2026}, howpublished={Preprint}, url={https://philpapers.org/rec/BENTCM-7}}
@book{jeffrey1983logic,
  author    = {Jeffrey, Richard C.},
  title     = {The Logic of Decision},
  edition   = {2nd},
  publisher = {University of Chicago Press},
  year      = {1983}
}

@misc{Aethus, author = {Benander, Ajax}, title = {Aethic Reasoning: A Comprehensive Solution to the Quantum Measurement Problem}, howpublished = {Preprint}, year = {2026}, url = {https://philpapers.org/rec/BENARA-5}, note = {The full development; PhilPapers record BENARA-5}}

@article{elga2000selflocating,
      author = {Elga, Adam},
      title = {Self-locating belief and the {S}leeping {B}eauty problem},
      journal = {Analysis}, volume = {60}, number = {2}, pages = {143--147}, year = {2000}}

@incollection{gibbard1978counterfactual, author={Gibbard, Allan and Harper, William L.}, title={Counterfactuals and Two Kinds of Expected Utility}, booktitle={Foundations and Applications of Decision Theory}, publisher={Reidel}, year={1978}}

@book{joyce1999foundations, author={Joyce, James M.}, title={The Foundations of Causal Decision Theory}, publisher={Cambridge University Press}, year={1999}}

@article{lewis1981causal, author={Lewis, David}, title={Causal Decision Theory}, journal={Australasian Journal of Philosophy}, volume={59}, pages={5--30}, year={1981}}

@incollection{nozick1969newcomb, author={Nozick, Robert}, title={Newcomb's Problem and Two Principles of Choice}, booktitle={Essays in Honor of Carl G. Hempel}, editor={Rescher, Nicholas}, publisher={Reidel}, year={1969}}

@misc{soares2017fdt, author={Yudkowsky, Eliezer and Soares, Nate}, title={Functional Decision Theory: A New Theory of Instrumental Rationality}, year={2017}, note={arXiv:1710.05060}}

@article{wolpert2011lesson, author={Wolpert, David H. and Benford, Gregory}, title={The lesson of Newcomb's paradox}, journal={Synthese}, volume={190}, number={9}, pages={1637--1646}, year={2013}, note={per deposited reference list (referee-transcribed), eyeball at merge}}

@article{wood_spekkens_2015, author={Wood, Christopher J. and Spekkens, Robert W.}, title={The lesson of causal discovery algorithms for quantum correlations: causal explanations of Bell-inequality violations require fine-tuning}, journal={New Journal of Physics}, volume={17}, pages={033002}, year={2015}}

@misc{yudkowsky2010tdt, author={Yudkowsky, Eliezer}, title={Timeless Decision Theory}, year={2010}, note={Machine Intelligence Research Institute, intelligence.org/files/TDT.pdf; per deposited reference list (referee-transcribed)}}

@misc{GoldenScience2022,
  author       = {Benander, Ajax},
  title        = {Golden Science},
  year         = {2022},
  month        = {November},
  howpublished = {Self-hosted manuscript},
  url          = {https://aethus.org/aethic-code/history/GoldenScience_Nov2022.pdf},
  note         = {Title page August 2022, the date of the first version of the source; this version created November 2022. Not deposited to a
                  preprint server; cited as a dated developmental record rather than
                  as a source of results.}
}
% R891 (Fable 5.1, Oct 30 2026), at the author's direction: aethus.org link updated to the file's new location. Previous: https://aethus.org/aethic-code/files/GoldenScience_Nov2022.pdf
\end{filecontents}

\documentclass[11pt]{article}
\usepackage[margin=1.05in]{geometry}
\usepackage{amsmath, amssymb, amsthm}
\newtheorem{principle}{Principle}
\newtheorem{definition}{Definition}
\newtheorem{lemma}{Lemma}
\newtheorem{proposition}{Proposition}
\newtheorem{theorem}{Theorem}
\newtheorem{corollary}{Corollary}
\usepackage[hidelinks]{hyperref}
\usepackage{microtype}
\usepackage{enumitem}
\usepackage{booktabs}
\usepackage{array}
\usepackage[style=numeric, sorting=none]{biblatex}
\addbibresource{refs.bib}
\addbibresource{dec_refs.bib} % local standalone coverage; refs.bib above wins collisions

\title{Newcomb's Problem Without a Fixed Past:\\ \large A Decision Rule From Observer-Indexed Determinacy}
\author{Ajax Benander\thanks{Email: \href{mailto:akb9408@rit.edu}{akb9408@rit.edu}, Website: \url{https://aethus.org}. This paper extracts, as a self-contained proposal in decision theory, one application of a larger framework \cite{AethusBrief, Aethus}; the framework's quantum-mechanical claims play no role in the argument below and are cited only for the record. The anthropic apparatus of the companion \cite{Nexus} enters as one of two routes to the same verdict and is likewise not presupposed.}}
\date{First drafted July 2026\\ {\small This version: October 2026 --- working draft}}

\begin{document}
\maketitle

\begin{abstract}
\noindent Newcomb-type problems require an account of which dependencies remain in force when a candidate act is evaluated. This paper proposes a settling rule motivated by observer-indexed determinacy, under which an unresolved outcome remains open relative to an information state regardless of its temporal location. On finite acyclic factor models with specified dependency roles and binding modes, the rule fixes the candidate act, retains binding couplings, jointly averages nonbinding couplings, and conditions on stated evidence. I characterize when the resulting outcome distributions are invariant across acts and hence support a common-distribution dominance argument. Under explicit assumptions, the rule recommends one-boxing in opaque Newcomb, smoking in the lesion case, and two-boxing in transparent Newcomb with positive prediction error. The transparent case exposes a conflict between current-state act evaluation and policies with higher ex-ante expected utility. An analysis of Egan's psychopath-button problem distinguishes evaluating each act at its own settled state from standard ratifiability; the result depends on a stated endpoint trait law. Observer-indexed determinacy motivates the proposal, while the binding assignments and normative justification remain additional commitments.
\end{abstract}

\section{The disagreement and the shared premise}
\label{sec:disagreement}

Newcomb's problem, in Nozick's original presentation \cite{nozick1969newcomb}, sets a near-perfect predictor against a single agent. The predictor---traditionally a supercomputer with an essentially flawless record---has already placed, in an opaque box (Box B), either one million dollars or nothing, according to its prior assessment of what the agent will do. A second, transparent box (Box A) contains a thousand dollars, visibly. If the predictor expected the agent to take only Box B, it placed the million; if it expected the agent to take both boxes, it left Box B empty. The agent now chooses: only Box B (\emph{one-boxing}), or both boxes (\emph{two-boxing}).

Two lines of reasoning pull in opposite directions, and each is locally compelling. Expected-utility reasoning observes that one-boxers, across the predictor's record, reliably walk away with a million and two-boxers reliably walk away with a thousand; so one-box. Dominance reasoning observes that the contents of Box B are already set, whatever they are, so that taking Box A as well yields a guaranteed extra thousand in every state of the world; so two-box. Formalized, these become the two standing decision theories, and both maximize an expected utility; they differ in the probabilities it is computed with. Evidential decision theory, in the tradition of Jeffrey \cite{jeffrey1983logic} and more recently Ahmed \cite{ahmed2014evidence}, computes it with the act-conditional probabilities and one-boxes, because choosing only Box B is evidence that the predictor foresaw one-boxing and loaded the million. Causal decision theory, originating with Gibbard and Harper \cite{gibbard1978counterfactual} and developed by Lewis \cite{lewis1981causal} and Joyce \cite{joyce1999foundations}, computes it with probabilities that hold the causally independent facts fixed, which in this problem returns the dominance verdict: two-box, because the prediction is causally sealed in the past and no present act can reach back to alter it.

The disagreement has been remarkably durable \cite{wolpert2011lesson}, and its durability is the datum this paper takes seriously. When two carefully formalized theories persist in opposite verdicts for fifty years, the productive suspicion is not that one side has simply failed to see the other's argument, but that both share a premise doing more work than either notices. Here that premise is explicit and common to both: \emph{the prediction, having been made in the past, is a determinate fact at the moment of choice.} Causal decision theory affirms it directly---the fixed past is its dominance argument's engine. Evidential decision theory does not deny it either; it accepts that the prediction is settled and argues only about what the agent's act reveals concerning that settled fact. The two camps agree that Box B's contents are, right now, one definite way; they disagree only about whether the agent's choice bears news about which way. This paper's proposal is to examine the shared premise itself. Two things are kept apart from the start: what the ontology permits once the premise is dropped, and which evaluation policy the paper adopts. Removing the premise motivates a decision rule without selecting one, the rule is stated as an operator on finite acyclic factor models (Section~\ref{sec:operator}), and its central normative burden, why an agent should adopt it when a competing policy does better ex ante in the transparent case, is taken up in Section~\ref{sec:defense}.

\section{The proposal: determinacy without tense}
\label{sec:proposal}

\noindent \textit{What kind of result this is.} Nothing here shows empirically that the prediction was not already determinate. Because the sealed box contains records, the reading on which its contents are open and the reading on which they are determinate-but-unread are indistinguishable by any observation the problem permits; the two-boxer who retains tensed determinacy has not been refuted. What is offered is a coherent alternative on which ``past but unread'' does not entail ``determinate for this agent,'' together with a decision rule that alternative motivates and the verdicts the rule returns across a family of cases. The result is a reconceptualization, and the argument for it is the usual one for a framework: that it discharges the puzzle without ad hoc machinery and sorts the neighboring cases correctly.

\noindent \textit{Three questions, kept apart.} The paper answers three questions, and the answer to the first does not answer the other two. What does the ontology permit or motivate? Tense-free blankness removes the tensed ground of the premise both standard theories share, that the prediction is fixed for the deciding agent because it was made before the act (this section). What settling policy does the paper adopt? Evaluation of each act at the information state its own settling induces, with stated couplings retained or detached by their mode of binding (Definition~\ref{def:settle}, Principle~\ref{prin:rule}). Why should an agent use that policy? That is a normative question, answered by a normative argument with its costs stated (Section~\ref{sec:defense}), and not by the ontology.

\noindent \textit{Which indexing the paper needs.} Four claims are easily run together under ``observer-indexed determinacy,'' and they differ in strength. \emph{Indexed representation}: the agent's information state represents the box as blank; any partial-information model has it. \emph{Indexed truth}: whether ``the box contains a million'' is true is relative to a state. \emph{Indexed determinacy}: whether the box's contents are \emph{determinate} is relative to a state, with no further fact of the matter absent an index. \emph{Objective determinacy}: there is a state-independent determinate fact about the contents, whatever any state says. The paper requires the third and denies the fourth. A reader who holds the fourth is not thereby committed to two-boxing, nor a reader who holds the third to one-boxing: the verdict turns on the settling rule, not on the reading of ``blank'' alone. Write the act-evaluation law as $Q_a(y) = P_A(Y = y \mid \mathrm{Settle}(A, a))$. One settling rule preserves the prediction's marginal across acts and supports two-boxing; another preserves the predictor's branch-conditional accuracy and supports one-boxing; treating the prediction as blank selects neither. The thesis is accordingly in three parts. Tense-free determinacy motivates evaluating acts through indexed settlement. A specified settlement policy, the branch-conditional one with its binding assignments, determines which dependencies survive. The verdicts follow from that policy and the supplied dependence structure. What the indexing removes is the tensed ground of marginal preservation, that the prediction is determinate independently of the act because it was made before the act; its untensed ground, that an evaluative intervention leaves the prediction's mechanism untouched, is not removed by the indexing and is met as a modeling commitment in Section~\ref{sec:defense} (the typicality reading of the same placement is Appendix~\ref{sec:anthropic}). Indeterminacy is not claimed to entail the rule.

The shared premise says that a past occurrence is a present determinate fact, and it draws its force from an assumption so ordinary it is rarely stated: that determinacy attaches to positions in time, so that once an event has happened its outcome is settled for every later moment, full stop. Call this the \emph{tensed-determinacy} assumption. The proposal is that it can be replaced, coherently and with independent motivation from the framework this paper draws on \cite{Aethus}, by \emph{observer-indexed} determinacy: an outcome is determinate relative to an information state\footnote{A naming note, recorded once: \emph{information} here is not to be confused with the quantitative, Shannon-theoretic kind --- the states are logical objects, individuated by stated content and deductive closure, with truth relativized to them. It is to establish exactly this that the fuller development gives the states a proper name, \emph{Aethus}, plural \emph{Aethae} \cite{Aethus}, and the name is used below wherever the index changes and the change is the point: the agent's deliberative Aethus, the successor Aethus after settling, the spectator's Aethus against the agent's.}, the total body of what a given observer's situation has settled, with everything unsettled left open. The apparatus is minimal and is developed as a self-contained semantics in the companion \cite{AethusSemantics}. An information state assigns to each attribute either a definite value (the attribute is \emph{stated}) or no value (it is \emph{blank}); a blank attribute is not a hidden definite value the observer happens not to know but is unsettled relative to that state, standing in superposition over its possible values until something settles it; and \emph{blanks lack tense}: whether the referent of a blank sits in the observer's past, present or future makes no difference to its being blank.
% >>> ADDED (Fable 5, Aug 2026, author-directed): the long manuscript's
% compact formulation of tense-free blankness carried verbatim-in-substance.
% AETHIC-PROV: author-content (the line is the author's own), ai-placed.
The comprehensive development compresses the doctrine into a single clause worth carrying whole: \emph{blankness-resolution, not pastness, is the Markovian decider of determinacy} \cite{Aethus}. Determinacy is decided at the resolving of the blank --- wherever in time the resolved event sits --- and nothing upstream of that resolution has pre-decided it.

Applied to Newcomb: the predictor's placement sits in the agent's proper-time past, the box is sealed, and the agent has not seen its contents, so relative to the agent's state at the moment of decision the prediction's content is a blank attribute, in superposition over ``million'' and ``nothing.'' The tensed-determinacy assumption says otherwise only because it converts pastness into settledness, and that conversion is the move the proposal declines. What settles the blank is the agent's decision, through closure and weight. An information state contains everything its records entail; with the predictor's near-perfect record stated and, once the agent decides, the agent's own choice stated, a fully resolved chooser occupies a state whose residue places weight $1 - \varepsilon$ on the matching prediction and $\varepsilon$ on the mismatch, $\varepsilon$ the predictor's error rate.\footnote{With act-dependent error rates --- $\varepsilon_1$ given one-boxing, $\varepsilon_2$ given two-boxing --- the one-boxing verdict requires only $\varepsilon_1 + \varepsilon_2 < 1 - t/M$, about $0.999$ at the canonical payoffs: only the \emph{sum} of the error rates matters.} Near-perfection is not entailment, so what closure delivers is weight concentration, not statedness. At the perfect endpoint the mismatch configurations are excluded as jointly invalid, a declared constraint of the model and not the limit of the positive-error weightings: the distributions converge as $\varepsilon \to 0$, but convergence of a weighting does not by itself create a hard exclusion, which is a separately stated feature of the model (Definition~\ref{def:settle}, Corollary~\ref{cor:transparent-sort}). Partial resolve is the same mechanism at interior points: the residue's weights track the firmness of the commitment, so that under the affine resolve assumption below the proposal delivers a \emph{resolve-gradient}, a continuous family of verdicts between open deliberation and full commitment, whose mathematics the deliberational-dynamics accounts share and whose reading as a weighting over status-types is this paper's.

\begin{principle}[The decision rule]\label{prin:rule}
For a deliberative state $A$ and a candidate act $a$, let $Q_a$ be the settled law of the payoff-relevant attributes $Y$ returned by the settling operator of Definition~\ref{def:settle}, which fixes the act, retains the couplings that bind in the evaluation mode, detaches the others jointly, and conditions on the stated evidence; $a$ is \emph{admissible}, $a \in \mathcal{A}_{\mathrm{adm}}(A)$, when its settling has a nonempty residue of positive total weight. The rule is admissible-act maximization of the settled expected utility,
\[
D(A) \;=\; \arg\max_{a \in \mathcal{A}_{\mathrm{adm}}(A)} \sum_{y} Q_a(y)\, u(a, y),
\]
each act evaluated at the state its own settling induces, with the stated correlations binding on the settling per their mode (Definition~\ref{def:binding}). An empty residue and a nonempty residue of zero weight alike supply no normalized expected utility without a further convention, and the paper's convention is that such an act is inadmissible.
\end{principle}

\noindent One-boxing follows at once, and the arithmetic is displayed because it earns three results in half a page. Write $t$ for the transparent box's payoff, $M$ for the opaque prize, $\varepsilon$ for the predictor's error rate. With the accuracy coupling retained, the resolved one-boxer's residue yields $\mathrm{EU}(\text{one-box}) = (1-\varepsilon)M$ and the resolved two-boxer's $\mathrm{EU}(\text{two-box}) = t + \varepsilon M$, so one-boxing wins exactly when
\begin{equation}
\varepsilon \;<\; \tfrac{1}{2} - \tfrac{t}{2M},
\end{equation}
which for the canonical payoffs ($t = 10^3$, $M = 10^6$) is $\varepsilon < 0.4995$: any predictor better than chance by more than a hair suffices. The resolve-gradient is the same display continued. For resolve $r \in [0,1]$ toward one-boxing---the weight the agent's state places on the one-box value of the act attribute, a state-indexed quantity and neither a mixed-strategy probability nor a credence in an external fact---the stated accuracy coupling gives the prediction marginal $\varepsilon + r(1-2\varepsilon)$. The affine form is an assumption about the coupling: it is what results if the predictor's accuracy is read off the act attribute's weight linearly, and any other monotone response would preserve the endpoint verdicts while curving the interior. The expected utility is
\begin{equation}
\mathrm{EU}(r) \;=\; t + \varepsilon M + r\left[(1-2\varepsilon)M - t\right],
\end{equation}
affine in $r$, its slope changing sign at the threshold above, so the gradient and the threshold are one result, and since $\mathrm{EU}(r)$ is linear its optimum sits at an endpoint whenever the slope is nonzero: the rule never counsels deliberate half-heartedness. What the gradient adds over the mixed strategies the standard theories already contain is the semantic status of $r$, a feature of what the agent's state states about the act rather than of what the agent will do; deliberation as an evolving state converging to fixed points is Skyrms's deliberational dynamics \cite{skyrms1990dynamics}, and what the present treatment adds is an ontology for the state, not a new dynamics. Tensed determinacy can model graded deliberation too, as graded belief about a fixed but unknown past; the mathematics is shared, and the difference is in what $r$ represents.

This is not evidential decision theory's news value: the agent is not reasoning ``my choosing one box is good evidence that the million is there,'' but holding that, relative to the agent's own state, the prediction was never a settled fact about which the act could carry news. And it is not backward causation: there is no causal arrow from the present act to the past placement, because there is no single tenseless fact of the placement for such an arrow to target. The act--prediction coupling is not derived by the framework; it is encoded in the state as a joint-invalidity constraint, the mismatched configurations receiving small positive weight because the vignette stipulates a near-perfect predictor, and what the operator does is fix the act, retain the binding accuracy coupling, and compare the settled laws. Every one-boxing theory produces an expected-utility calculation once a coupling is supplied; the contribution here is the ontology on which the coupling binds the act without being a causal arrow, and the criterion of Section~\ref{sec:operator} that says when such a coupling binds and when it does not.

% >>> ADDED (Fable 5.1, Sep 2 2026, discipline pass): the positive principle behind the two disclaimers,
% stated once --- the correlation accounted for by joint invalidity, no direction and no mechanism --- so
% that ``not news, not backward causation'' reads as a consequence rather than a defense.
\begin{principle}[Joint coupling in place of causation]\label{prin:joint_invalidity}
The predictor's accuracy is not a causal arrow in either direction but a fact about which joint configurations of the agent's information state are valid: perfect prediction excludes the mismatch combinations, \emph{one-box and predicted two-boxing} and \emph{two-box and predicted one-boxing}, as jointly invalid relative to that state, while imperfect prediction retains them at the stipulated error weights, invalidity and low weight being distinct, and the act's settling selects among the survivors. A correlation is thereby accounted for by the \emph{joint coupling} of the combinations it constrains---the weights the state assigns across them---with no direction and no mechanism in the account. ``$A$ causes $B$'' is not rendered here as ``$\{A, \neg B\}$ is jointly invalid,'' which is a support implication that two effects of a common cause satisfy without either causing the other; the causality companion's distinctions, factor graph, exclusion structure and mechanism history, are preserved \cite{AethusCausality}. A correlation may constrain every combination without excluding any.
\end{principle}
\noindent Evidential decision theory reads the correlation as news; causal decision theory asks whether the act is a cause of the prediction, to which the answer is no. The present account takes neither reading: the correlation is a joint coupling, and what is stated about it is whether it binds under the mode by which the act is produced, the direction in play being the roles', supplied with the model. The same principle is the framework's account of correlation wherever it appears \cite{Aethus}, and its appearance here is the decision-theoretic face of a general replacement of causal structure by consistency structure, not a device invented for Newcomb.

\section{Model inputs and the settling operator}
\label{sec:operator}

\noindent The rule's inputs are a model, an evaluation mode, the stated evidence and a candidate act, and the operator is defined on finite acyclic factor models. Its verdicts are conditional on two kinds of structural information the model supplies and the theory does not derive: the roles, which attribute is produced from which, and the binding modes of the couplings with the act. The need for the second is not decorative. A theory that always distributes the state into its pockets before evaluating runs the causalist's dominance argument uniformly and two-boxes Newcomb; a theory that always compares whole residues abstains on the smoking lesion; each mode is right on one case, and the binding modes are the sorter that says which governs when.

\begin{definition}[Binding modes; act-robustness]
\label{def:binding}
A stated correlation involving the act carries a mode set $m(e) \subseteq \mathcal{M}$, the modes of act-production under which it is stated to hold, and it \emph{binds on the agent's settling} in mode $\mu$ iff $\mu \in m(e)$. A correlation stated as robust to every mode, $m(e) = \mathcal{M}$, holding whether the act issues from deliberation, habit, or a randomizing device, is the \emph{act-robust} special case; binding in the deliberative mode requires only $\mu_{\mathrm{del}} \in m(e)$. A correlation that does not bind in the evaluation mode binds only through its mediating attributes. The assumption that the predictor's near-certainty holds even against a randomizing agent, examined in Appendix~\ref{sec:randomization}, is exactly the statement that Newcomb's accuracy entry is act-robust; the present definition is that assumption generalized from coins to deliberation, and it is the substantiated form of what the functional-decision-theory literature calls subjunctive dependence \cite{soares2017fdt}: an act-robust stated coupling is a dependence stipulated to survive every mode of act-production, a stronger condition than functional decision theory's, which does not require invariance under arbitrary replacement of the decision-producing process; opaque one-boxing needs only that the branch-conditional accuracy law survive the deliberative evaluation being performed, and the all-modes condition is carried as a stipulation stronger than the verdict needs.
\end{definition}
\noindent Formally: for each primitive coupling $e$ the specification supplies a mode set $m(e) \subseteq \mathcal{M} = \{\text{deliberation}, \text{habit}, \text{randomizer}, \ldots\}$, the modes of act-production under which $e$ binds; $e$ is \emph{act-robust} iff $m(e) = \mathcal{M}$; and a settling operation performed in mode $\mu$ retains $e$ iff $\mu \in m(e)$. That retention rule governs the couplings with the act among their inputs: a coupling is retained exactly when the evaluation mode lies in its mode set, and the nonbinding ones are replaced jointly by the kernel from the act's production inputs, while every production entry of the act is replaced, as the definition below states. The theory does \emph{not} derive $m(e)$. It is structural information in the problem specification, as the edges of a causal model are, and every verdict below is conditional on it.


One sentence should stand beside the definition, because it answers the objection the definition most invites --- that act-robustness is causal dependence renamed, the causal arrow rebuilt and relabeled. Newcomb's accuracy entry is act-robust \emph{without} any act-to-prediction causation: that is the case where robustness and causation come apart, and it is why the present criterion one-boxes where the standard causal decision theories two-box (arrow-based frameworks that one-box exist --- Spohn's reflexive decision graphs the nearest --- but do so by enriching the arrow apparatus rather than by an invariance condition). Robustness concerns which stated couplings survive which settlings, not arrows; the Newcomb entry is the existence proof of the difference.

\noindent\emph{The settling operator.} Definition~\ref{def:settle} states the operator; three points about it are not visible in the definition. The roles are structural and not read off the joint law, which admits both $P(C)\,P(B \mid C)$ and $P(B)\,P(C \mid B)$ for a choice $C$ and contents $B$: with the roles assigning $C$'s production and a $C \to B$ coupling, settling preserves $P(B \mid C = c)$, the branch-conditional rule, and with $B$'s production and a $B \to C$ coupling it preserves $P(B)$, the marginal-preserving rule. The roles and mode sets are supplied with the model, not derived, and the paper does not read the roles as causal production. Deleting a non-binding coupling is not an option: for an act $C$ with $P(C = 1) = 0.1$ and $P(Y = C \mid C) = 0.9$, so that $P(Y = 1) = 0.18$, deletion leaves $Y$ uniform, while detachment keeps $P(Y = 1) = 0.18$.

\begin{definition}[The settling rule]\label{def:settle}
\emph{Model class.} A finite directed acyclic graphical model over the act $C$, the payoff-relevant attributes $Y$ and auxiliary attributes, in which each attribute has exactly one kernel given its parents, its inputs, so that the joint law is the product of these kernels and its factorization is fixed by the supplied graph and kernels, a joint law alone not determining it. The act's kernel $P_0(c \mid \mathrm{pa}(C))$ is its \emph{production entry}; a \emph{coupling} is the kernel $f_e$ of an attribute having $C$ among its parents, and carries a supplied mode set $m(e)$.
\emph{Inputs.} The model; the evaluation mode $\mu$, deliberation for this paper's agent; the stated evidence $E$, a set of attribute values; and a candidate act $a$.
\emph{Output.} The settled law $Q_a$ of $Y$, or the verdict that $a$ is inadmissible under the model.
\emph{Construction.} Every replacement kernel is computed from the original model before any kernel is replaced, and the three replacements are made simultaneously. (1) \emph{Fix}: the production entry is replaced by the point mass at $a$. (2) \emph{Retain}: every coupling with $\mu \in m(e)$ is kept, with its dependence on the act. (3) \emph{Detach}: let $O$ be the outputs of the couplings with $\mu \notin m(e)$ and $Z$ their parents outside $O \cup \{C\}$. These couplings are replaced jointly by one kernel of $O$ given the act's parents and $Z$, computed from the original kernels with the act integrated out under the original production entry, $K(o \mid \mathrm{pa}(C), z) = \sum_c P_0(c \mid \mathrm{pa}(C)) \prod_{e} f_e\bigl(o_e \mid c,\, o_{\mathrm{pa}(e) \cap O},\, z_{\mathrm{pa}(e) \cap Z}\bigr)$, so that dependencies among the variables of $O$ stay inside the block and only $Z$ and the act's parents are conditioning arguments. The product of the resulting kernels is then restricted to $E$ and marginalized to $Y$, giving $\widetilde Q_a$, and $Q_a := \widetilde Q_a / \sum_y \widetilde Q_a(y)$.
\emph{Failure.} If the normalizer is zero, the settled act has probability zero given the retained kernels and the evidence, and $a$ is \emph{inadmissible under the model}. That notion is model-relative: it coincides with invalidity under the third postulate when the model's zero-probability joint configurations are exactly its declared joint invalidities, as at the perfect-predictor endpoint, where $B = a$ is a declared constraint, and in general a zero-probability configuration need not be a formal contradiction. Outside the model class, in cyclic or feedback structures, the rule is not defined here, and no general algorithm for them is claimed.
\end{definition}

\begin{proposition}[Existence and uniqueness of the settled law]\label{prop:settle-exists}
On the model class of Definition~\ref{def:settle}, for each candidate act $a$ either $a$ is inadmissible under the model or $Q_a$ exists and is unique; and any sequential implementation of the three replacements gives the same $Q_a$, provided each replacement kernel is computed from the original model.
\end{proposition}
\begin{proof}
Each replacement kernel is a function of the original model alone: the point mass at $a$, the retained couplings themselves, and the detached kernel, a finite sum over the act's values of products of original kernels. The replacements act on disjoint sets of kernels, the production entry, the binding couplings and the non-binding ones, each attribute having exactly one kernel, so the replaced product is determined, and sequential implementations agree because none reads a kernel another has replaced. Conditioning on $E$ and marginalizing to $Y$ are functions of that product, whose normalizer is finite in a finite model, so the normalization exists exactly when the normalizer is positive. The freezing is needed: with $C$ fair and one non-binding coupling $Y := C$, detaching with the original production entry leaves $Y$ fair, whereas detaching after the production entry had been replaced by the point mass at $a = 1$ would give $Y = 1$.
\end{proof}

\section{The auxiliary-variable representation and across-act invariance}
\label{sec:invariance}

\begin{proposition}[Auxiliary-variable representation]\label{prop:aux}
On the model class of Definition~\ref{def:settle}, $Q_a$ is obtained as follows. Add to the model one copy $C^{\circ}$ of the act with the original production entry $P_0(c^{\circ} \mid \mathrm{pa}(C))$; fix $C = a$; let every coupling with $\mu \in m(e)$ read $C$ and every coupling with $\mu \notin m(e)$ read the same $C^{\circ}$, their kernels otherwise unchanged; leave every other kernel unchanged; incorporate the stated evidence; and marginalize $C^{\circ}$. Marginalizing the shared copy yields exactly the joint-detachment kernel $K$ of step (3), and separate copies, one per detached coupling, would not: they would make the detached outputs independent given their other parents, discarding the dependence the shared act induced among them (Appendix~\ref{sec:withdrawn-operators}). If the act is instead given an independent prior of full support in the augmented graph $G^{\circ}$, then $Q_a(Y) = \widehat P(Y \mid C = a, E)$ for every admissible $a$, $\widehat P$ the augmented law.
\end{proposition}
\begin{proof}
In the augmented joint the factor of the detached outputs is $\prod_e f_e(o_e \mid c^{\circ}, o_{\mathrm{pa}(e) \cap O}, z_{\mathrm{pa}(e) \cap Z})$ with $C^{\circ} \sim P_0(\cdot \mid \mathrm{pa}(C))$, and summing over $c^{\circ}$ gives the kernel $K(o \mid \mathrm{pa}(C), z)$ of step (3); the retained factors read $C = a$ as step (2) has them, and the fix is step (1). With an independent full-support prior $\rho$ on $C$ the augmented joint is $\rho(c)$ times the settled joint at $c$, so conditioning on $C = a$ and $E$ returns the settled law at $a$, the prior cancelling.
\end{proof}
\begin{corollary}[A graphical sufficient condition for invariance]\label{cor:dsep}
In the augmented graph $G^{\circ}$ with the act exogenous, if the payoff-relevant attributes $Y$ are d-separated from $C$ given the stated evidence $E$, then $Q_a(Y)$ is the same for every admissible act and the hypothesis of Lemma~\ref{lem:invariance} holds. The condition is sufficient and not necessary, cancellation along open paths being possible without faithfulness; it is the paper's graphical criterion for invariance, and the remark after Corollary~\ref{cor:dominance-transfer} records why no simpler reading replaces it.
\end{corollary}
\begin{proof}
d-separation in $G^{\circ}$ gives $\widehat P(Y \mid C = a, E) = \widehat P(Y \mid E)$, which does not depend on $a$, and by Proposition~\ref{prop:aux} the left side is $Q_a(Y)$.
\end{proof}

\begin{lemma}[Across-act invariance from factorization]
\label{lem:invariance}
Let $Y = (Y_1, \ldots, Y_n)$ collect the payoff-relevant blank attributes, and for each admissible act $a$ let $\widetilde Q_a(y)$ be the unnormalized law of $Y$ after the act is settled in the deliberative mode by the operator of Definition~\ref{def:settle}, the stated evidence is incorporated, and the auxiliary attributes are marginalized. If
\[
\widetilde Q_a(y) = k_a\, h(y) \quad \text{for every admissible } a, \qquad k_a > 0, \quad 0 < \sum_y h(y) < \infty,
\]
with one and the same $h$ for every act, then the settled law of $Y$ is the same across candidate acts: \[ Q_a(y) = \frac{h(y)}{\sum_{y'} h(y')} \quad \text{for every admissible } a. \] Conversely, if the settled laws coincide, $Q_a = Q$ for every admissible $a$, and every normalizer $\sum_y \widetilde Q_a(y)$ is positive and finite, then $\widetilde Q_a = k_a h$ with $h = Q$ and $k_a = \sum_y \widetilde Q_a(y)$. Proportionality of the unnormalized settled laws and equality of the settled laws are therefore equivalent. \end{lemma} \begin{proof} Normalizing $\widetilde Q_a$ divides by $k_a \sum_{y'} h(y')$, and $k_a$ cancels. For the converse, $\widetilde Q_a(y) = \bigl(\sum_{y'} \widetilde Q_a(y')\bigr)\, Q_a(y) = k_a\, Q(y)$. \end{proof}

\begin{corollary}[Dominance transfer]
\label{cor:dominance-transfer}
Under the hypothesis of Lemma~\ref{lem:invariance}, with $Q$ the common settled law, if $u(a, y) \ge u(b, y)$ for every $y$ then $\sum_y Q(y)\, u(a, y) \ge \sum_y Q(y)\, u(b, y)$, with strict inequality whenever the pointwise inequality is strict on a set of positive $Q$-mass: pointwise dominance transfers to expected utility.
\end{corollary}

\noindent\emph{Why statedness does not block a path.} An earlier draft proposed a path criterion on which a dependence of $Y$ on the act is blocked whenever it runs through a stated attribute. That reading is false in general, and the d-separation condition is what replaces it. Let $C$ and $Y$ be independent fair bits and let a retained binding coupling $S$ read both, with $P(S = 1 \mid C, Y) = 0.9$ when $C = Y$ and $0.1$ when $C \ne Y$, and let $S = 1$ be stated. Settling $C$ gives $Q_0(Y = 1) = 0.1$ and $Q_1(Y = 1) = 0.9$, both normalizers equal to $\tfrac12$: the settled laws differ although the only attribute between the act and $Y$ is stated, because conditioning on the common output opens the path $C \to S \leftarrow Y$. In $G^{\circ}$ the collider is open given $S$, so Corollary~\ref{cor:dsep} correctly declines to certify invariance; had $S$ been nonbinding it would read $C^{\circ}$, the path would be closed, and the settled laws would coincide at $\tfrac12$, as the accompanying script checks. A stated output whose kernel reads both the act and $Y$ therefore leaves a likelihood in both, statedness of a binding output does not supply invariance by itself, and preservation of the original law, $Q_a(Y) = P_A(Y)$, is not claimed at all (Appendix~\ref{sec:withdrawn-operators}). Whether a graphical criterion broader than d-separation in $G^{\circ}$ characterizes invariance on this model class is left as a research question; it is not used as a premise anywhere below, where each benchmark's factorization is verified or its settled laws are calculated directly.

\begin{proposition}[Dominance recovery]
\label{prop:dominance}
Invariance of the \emph{joint} law over the payoff-relevant blanks across candidate acts, Lemma~\ref{lem:invariance}, is the condition under which pocket-wise dominance is sound uniformly and without further assumption: there, distributing the state into its pockets before evaluating loses nothing, and the pocket-wise dominant act is the rule's own verdict. Where only the marginals are invariant, dominance is sound for additively separable utility and not in general. The framework thereby states a sufficient domain for dominance, and its boundary.
\end{proposition}

\begin{theorem}[Decision classification]
\label{thm:classification}
Let the binding modes $m(e)$ be supplied, and let every invariance below mean the across-act invariance of Lemma~\ref{lem:invariance}, whose factorization hypothesis is what each clause consumes and is verified for each benchmark. Then, for the act's settling in the deliberative mode:
\begin{enumerate}
\item[(a)] If the factorization hypothesis holds for $Y$---verified model by model, and supplied by Corollary~\ref{cor:dsep} wherever $Y$ is d-separated from the act in $G^{\circ}$ given the evidence---settling is invariant across acts and pocket-wise dominance applies by Corollary~\ref{cor:dominance-transfer}.
\item[(b)] If the factorization fails, equivalently, by the converse in Lemma~\ref{lem:invariance}, if the settled laws differ between two admissible acts, the whole-residue evaluation of Principle~\ref{prin:rule} is required, and pocket-wise dominance is not guaranteed. The failure is established on the complete marginalized product $\widetilde Q_a$; a single retained factor's dependence on $y$ does not by itself establish it.
\item[(c)] A route to (a): if every apparent act-dependence is mediated through a stated attribute $M$ that screens it off---the settled evaluation family factorizing as $\widetilde Q_a(y, m) \propto f(a, m)\, g(y, m)$ for every admissible $a$, a statement about the post-settlement laws and not about the observational distribution, which statedness alone does not supply---conditioning on $M$ makes the surviving act-dependent contribution constant in $y$ and the factorization holds. Without faithfulness, cancellation along an open path can also leave the settled laws equal across acts, and whole-residue evaluation still decides.
\end{enumerate}
The benchmark verdicts fall into the two regimes, equal or unequal settled laws, screening being one sufficient route to equality. \emph{Newcomb}: the predictor's accuracy is stated act-robust, and the complete settled laws of the prize differ between the acts, $(0.99, 0.01)$ over full and empty under one-boxing against $(0.01, 0.99)$ under two-boxing, so proportionality fails, case (b), one-box. \emph{The lesion}: the correlation between act and outcome runs through the lesion, which stays blank, and the stated desire screens it off, case (c), smoke---provided the desire is a process variable $R$ for which act and lesion are independent given it, $C \perp L \mid R$; the common outcome law is then displayed directly, $Q_{\mathrm{smoke}}(L) = Q_{\mathrm{refrain}}(L) = P(L \mid R = r)$ with $r$ the stated resolve, the act adding nothing to what the stated resolve already says about the lesion. Statedness alone does not give that: with $P(L) = \tfrac12$, $P(A \mid L) = 0.9$ and $P(A \mid \neg L) = 0.1$, the resolve $r = \tfrac12$ can be stated while $P(L \mid A, r) = 0.9 \neq P(L \mid r) = \tfrac12$. \emph{Transparent Newcomb}: the prize is already stated, and with positive predictor error the settled laws of the remaining payoff-relevant blanks coincide across the acts given the stated prize, so proportionality holds, case (a), two-box, parting from the updateless, policy-level comparator's one-boxing for $0 < \varepsilon < (M - t)/(2M - t)$ (Proposition~\ref{prop:shared-model}); at the perfect endpoint the admissibility rule of Principle~\ref{prin:rule} governs, as Corollary~\ref{cor:transparent-sort} states.
\end{theorem}
\begin{proof}
(a) is Lemma~\ref{lem:invariance} with Corollary~\ref{cor:dominance-transfer} and Proposition~\ref{prop:dominance}; (b) is the failure of the lemma's factorization, exhibited directly for each benchmark; (c) is the factorization holding through a stated screening attribute. There are two regimes, equal and unequal settled laws, and screening is one sufficient route to the first, not a third case. The theorem is an algebraic classification once the evaluation laws have been obtained; Corollary~\ref{cor:dsep} is the graphical sufficient condition, and no broader graphical procedure is claimed. The verdicts are the regimes' instances.
\end{proof}

\section{Opaque Newcomb, the lesion, and transparent Newcomb}
\label{sec:cases}

Two further cases discriminate decision theories, and the operator is now run on both beside opaque Newcomb. In the smoking-lesion case, a gene causes cancer and, independently, a taste for smoking; smoking itself is causally inert; the agent enjoys smoking, and the question is whether to smoke. In transparent Newcomb, both boxes are open and their contents visible before the choice is made.

\begin{corollary}[Newcomb]
\label{cor:newcomb-sort}
The accuracy entry is stated act-robust --- the coin-inclusive strengthening of the vignette's predictor, examined in Appendix~\ref{sec:randomization} --- hence pocket-defining --- an entry whose robustness precludes act-independent pocketing of the target family: no prediction-pockets exist independently of the act, invariance fails for the payoff-relevant blank, the common-distribution dominance argument is unavailable, and residue comparison governs --- one-box.
\end{corollary}

\begin{corollary}[The lesion, in three variants]
\label{cor:lesion-sort}
The vignette stipulates a common cause and only a common cause: the lesion reaches the act solely through the desire, and nothing in the vignette states the desire--act propensity as robust to deliberate settling. That reading of the stipulations is a premise, defended in those two clauses. (i) \emph{Desire stated}: the lesion's weight is already updated on it, and in $G^{\circ}$ the lesion is d-separated from the act given the stated desire (Corollary~\ref{cor:dsep}), so invariance holds and the common-distribution dominance argument governs; smoke. (ii) \emph{Desire blank but introspectable}: the resolve-gradient of Section~\ref{sec:proposal} \emph{is} the tickle \cite{eells1982rational} --- settling-from-desire progressively states the desire, the leaning the gradient tracks being the desire manifesting --- so by settlement the case-(i) screen is in place; smoke. (iii) \emph{Influence bypassing introspection}: the agent's own resolve is a stated attribute of the agent's own state --- a mild luminosity premise, booked here explicitly: the resolve is the state's own weight profile over the act, so the state's carrying it is nearly analytic; and the premise does not reopen the bookability door the semantic companion closes \cite{AethusSemantics}, a resolve being a commitment attribute, the state's leaning, not a stored conditional verdict about its own weights, and it is verdict-type entries the triviality-blocking excludes --- and process-mediated influence reaches the act only through the process whose framework-representation the resolve is, so the resolve screens, in the metatickle and ratifiability lineage \cite{eells1982rational, jeffrey1983logic}, exactly when that routing holds as the independence $C \perp L \mid R$. The resolve's being stated is the luminosity premise; its screening is this further premise, and the second does not follow from the first, since a known resolve need not carry everything relevant to the lesion--act dependence, as the counterexample recorded with the classification theorem shows. Given both: smoke. An agent who cannot state even their own resolve is outside this framework's reach.
\end{corollary}

\begin{corollary}[Transparent Newcomb, with an admissibility rule] \label{cor:transparent-sort} The outcome is stated. With strictly positive predictor error and branch-specific support, every mismatch pocket carrying positive weight, which positive aggregate error does not by itself guarantee, a mismatch pocket survives every settling, both acts are admissible, the payoff-relevant attribute is itself stated so that its settled law is the point mass at the observed contents under both acts, invariance holds trivially, the common-distribution dominance argument governs, and the verdict is two-box. At the perfect-predictor endpoint the retained accuracy constraint is $B = a$, so settling two-boxing with the million visible yields the state $B = 1$, $a = 0$, $B = a$, jointly invalid, with empty residue and no expected utility; the paper's policy is that an act whose settling is jointly invalid with the retained entries is inadmissible, so at that endpoint one-boxing is the only admissible act, and the corollary's verdict changes with the case: two-box with positive error whatever is visible; at the perfect endpoint, one-box only with the million visible, and two-box only with the opaque box empty, where under the constraint $B = a$ on the actual act two-boxing is the sole admissible act. Under the shared model of Definition~\ref{def:shared-model} the variable the operator settles at the full box is the policy component $\pi_{\mathrm{full}}$: $\mathrm{Settle}(A, a)$ replaces its production entry by the point mass at $a$, retains the prediction's binding coupling to it, and conditions on the stated contents, so that the settled law of $B$ is the point mass at the observed contents under both acts and the comparison is $M$ against $M + t$; the empty-box act is $\pi_{\mathrm{empty}}$, which perfect accuracy leaves unconstrained: both empty-box acts are compatible with the target, and two-boxing is selected by payoff dominance, $t$ against nothing, not by admissibility; settling an act sets its policy component and leaves the other as the policy had it. The limit of the positive-error evaluations is not used to evaluate the endpoint, a conditional evaluation's limit being no evaluation at the point where its conditioning event becomes impossible. \end{corollary}

\noindent\emph{What each verdict consumes.} Among the rule's inputs three are distinct: aggregate historical accuracy, accuracy conditional on each act, and the robustness of those conditional laws under the settling that evaluates a candidate act. The second and third are premises, not consequences of the first. A predictor that always predicts two-boxing, among agents of whom one in a thousand one-boxes, is 99.9 percent accurate in aggregate while its opaque box is always empty, so that $V(1) = 0$ and $V(2) = t$. Table~\ref{tab:inputs} lists the essential additional input of each benchmark.
\begin{table}[ht]
\centering
\small
\caption{The essential additional input each benchmark verdict consumes, beyond the decision rule.}
\label{tab:inputs}
\begin{tabular}{p{0.24\textwidth} p{0.66\textwidth}}
\hline
Case & Essential additional input \\
\hline
Opaque Newcomb & Branch-conditional accuracy, retained under settling \\
Smoking lesion & A screening relation $C \perp L \mid R$, or a detached act-production mechanism \\
Transparent Newcomb & The visible contents retained while each candidate act is evaluated ; two-box with positive error, the admissibility rule governing the perfect endpoint \\
Egan's button & A resolve-dependent trait law, and jointness: an act opposite to its resolve is tagged , and the endpoint trait law (J2), the trait weight given the terminal resolve being $q_1$ whatever the prior lean \\
Randomized Newcomb & Accuracy extended to acts a randomizer produces \\
\hline
\end{tabular}
\end{table}


Three verdicts --- one-box, smoke, and two-box under positive predictor error, the perfect endpoint governed by the admissibility rule --- from one result: given the specified evaluation laws, their equality or inequality across acts determines whether the common-distribution dominance argument applies (Lemma~\ref{lem:invariance}), and the three verdicts are the two regimes' instances under the stated dependence annotations, with Corollary~\ref{cor:dsep} the graphical sufficient condition for the first regime. Three premises were stated along the way and are recorded here as premises: the non-robustness reading of the lesion vignette (two clauses, above); the resolve-luminosity premise, that the resolve is stated (mild, near-analytic, but a premise); and the screening premise, that the stated resolve is a sufficient screen, $C \perp L \mid R$, which does not follow from its being stated. The lemma is then near-analytic given the definition, and deliberately so: the substance of the sorting lives in those classifications, stated where substance belongs rather than hidden in a proof. One remark locates the criterion in a larger program. The sorting above consulted no causal production, though it consulted the supplied roles and mode sets, directed evaluative structure that is supplied rather than derived: it consulted which stated couplings survive which conditionings, and the classical triad of causal stories --- the act causes the outcome, the outcome causes the act, a common cause produces both --- did no work anywhere in it. That is the shape Wood and Spekkens argued Bell correlations force on causal explanation \cite{wood_spekkens_2015}, met here at the heart of classical decision theory: under the relocation of determinacy this paper runs on, the three modes are not extended but dissolved --- the companion's thesis \cite{Aethus}, not argued here --- and the same criterion that sorts these cases is what the dissolution leaves behind.

The screening identity the lesion instances twice, the lesion screening smoking from cancer and the stated desire screening the act from the lesion, is derived for factorized states in the causality companion, Reichenbach's certificate applied at the mediator being no news to the tickle literature \cite{reichenbach1956direction, eells1982rational, AethusCausality}; nothing here depends on the companion's further results.

Transparent Newcomb sharpens the anchor claim from the other direction. With both boxes visible and positive predictor error, the prediction's content is stated, nothing relevant is blank, the settling has nothing to select among, and the proposal returns the causalist verdict: take both. At the perfect endpoint the admissibility rule governs instead, one-boxing being the only admissible act with the million visible and two-boxing with the opaque box empty. The functional family divides here. The updateless, policy-level evaluation of Yudkowsky and Soares's presentation one-boxes on a full box \cite{soares2017fdt}, while the updateful variant characterized by MacDermott, Everitt and Belardinelli two-boxes, as this rule does \cite{macdermott2023characterising}; so the transparent case separates the rule from the updateless comparator and from nothing else in that family, and the comparison is made in one model, the shared transparent-predictor model below, with its own prediction target and an explicit error law, both evaluations computed in it.

\begin{definition}[Shared transparent-predictor model]\label{def:shared-model}
Box $A$ holds $t$ and is open; box $B$ holds $M > t$ or nothing and is transparent. The agent's \emph{policy} is a pair $(\pi_{\mathrm{full}}, \pi_{\mathrm{empty}})$, each in $\{\text{one-box}, \text{two-box}\}$. (i) \emph{Target:} the predictor predicts $\pi_{\mathrm{full}}$, what the agent would do upon seeing $B$ full. (ii) \emph{Error law:} the prediction equals $\pi_{\mathrm{full}}$ with probability $1 - \varepsilon$ and its opposite with probability $\varepsilon$, independently of everything else; $B$ is filled iff the prediction is one-box. (iii) \emph{Observation:} the agent sees $B$'s contents before acting. (iv) \emph{What each rule holds fixed:} the updateless comparator evaluates the policy pair, letting the prediction vary with $\pi_{\mathrm{full}}$ through (i)--(ii) and holding the error law fixed, the observation determining which component is varied; the settling rule of Principle~\ref{prin:rule} evaluates each act with $B$'s observed contents stated, under the admissibility rule at $\varepsilon = 0$. (v) \emph{Which variable each observation settles:} on a full box the operator settles $\pi_{\mathrm{full}}$, the predictor's target; on an empty box it settles $\pi_{\mathrm{empty}}$, and the predictor's target, $\pi_{\mathrm{full}}$, is left as the policy had it.
\end{definition}

\begin{proposition}[The two verdicts in the shared model]\label{prop:shared-model}
In the model of Definition~\ref{def:shared-model}: the updateless comparator sets $\pi_{\mathrm{full}} = $ one-box exactly when $(1 - \varepsilon) M + \varepsilon t > (1 - \varepsilon) t + \varepsilon (M + t)$, that is, when $\varepsilon < (M - t)/(2M - t)$, which for $M = 10^{6}$ and $t = 10^{3}$ is $\varepsilon < 0.4997$; and it sets $\pi_{\mathrm{empty}} = $ two-box always. The settling rule two-boxes on a full box for every $\varepsilon > 0$, the mismatch pocket then carrying positive weight, and one-boxes on a full box at $\varepsilon = 0$, by the admissibility rule; it two-boxes on an empty box always. The verdicts therefore agree at $\varepsilon = 0$ and on an empty box, and part on a full box for every $0 < \varepsilon < (M - t)/(2M - t)$, the updateless comparator one-boxing and the settling rule two-boxing. Ex ante, the one-box policy earns $(1 - \varepsilon) M + \varepsilon t$ and the two-box policy $(1 - \varepsilon) t + \varepsilon (M + t)$.
\end{proposition}
\begin{proof}
Under $\pi_{\mathrm{full}} = $ one-box the box is full with probability $1 - \varepsilon$, yielding $M$, and empty with probability $\varepsilon$, yielding $t$ under the dominant $\pi_{\mathrm{empty}}$; under $\pi_{\mathrm{full}} = $ two-box it is full with probability $\varepsilon$, yielding $M + t$, and empty with probability $1 - \varepsilon$, yielding $t$. The difference is $(1 - 2\varepsilon) M - (1 - \varepsilon) t$, positive iff $\varepsilon < (M - t)/(2M - t)$. For the settling rule, with the full box stated nothing payoff-relevant remains blank, so two-boxing yields $M + t$ against $M$ whenever it is admissible, and at $\varepsilon > 0$ the pocket in which the prediction was one-box and the act is two-box has weight $\varepsilon > 0$ and is admissible; at $\varepsilon = 0$ that pocket is invalid and one-boxing is the only admissible act. On an empty box both rules take $t$ over nothing.
\end{proof}

\noindent The two queries instantiate the operator as follows.
\begin{center}\small
\begin{tabular}{lll}
\toprule
Observation & Variable settled & Predictor's target \\
\midrule
Full box & $\pi_{\mathrm{full}}$ & $\pi_{\mathrm{full}}$ \\
Empty box & $\pi_{\mathrm{empty}}$ & $\pi_{\mathrm{full}}$, unchanged \\
\bottomrule
\end{tabular}
\end{center}
\noindent On a full box, $\mathrm{Settle}(A, a)$ replaces the production entry of $\pi_{\mathrm{full}}$ by the point mass at $a$, retains the prediction's binding coupling to $\pi_{\mathrm{full}}$, and conditions on $B$ full: the settled law of $B$ is the point mass at the observed contents under both acts whenever both are admissible, and the comparison is $M$ against $M + t$. On an empty box, $\mathrm{Settle}(A, a)$ replaces the production entry of $\pi_{\mathrm{empty}}$; the prediction's coupling reads $\pi_{\mathrm{full}}$, which the settling leaves as the policy had it, so the prediction's law is unchanged, both acts are admissible at every $\varepsilon$, and two-boxing is selected by payoff dominance, $t$ against nothing, not by admissibility. Perfect prediction constrains what the agent would do upon seeing a full box; it does not make two-boxing the only admissible empty-box act. This is why the model predicts the response to a specified observation rather than the act without qualification, the latter producing consistency problems in transparent variants. The rules differ in how observed content constrains counterfactual evaluation: the settling rule holds the observed contents fixed and settles the component the observation selects, whereas the updateless comparator varies $\pi_{\mathrm{full}}$ together with its predictive dependencies; neither returns the agent to an earlier information state, and the description below is of what each evaluation holds fixed. In the framework's terms the two evaluations are evaluations of two different Aethae, the one with the contents stated and the one in which the policy's effect on the prediction is open, but the difference is one of procedure, of what the counterfactual holds fixed, and not of information alone. The ex-ante figures record the cost of the settling rule's choice, since an agent of its kind is predicted to two-box on a full box and so seldom sees one: at $\varepsilon = 0.01$ the policy that one-boxes on full earns $990{,}010$ in expectation against $11{,}000$, the figures MacDermott, Everitt and Belardinelli report for the updateful and updateless variants, where precommitment is also discussed \cite{macdermott2023characterising}; the contribution here is the indexed interpretation and the specified operator, not the numerical tradeoff. The tradeoff is stated as the criterion's commitment and not concealed: the rule's full-box verdict changes between exact prediction, where one-boxing is the only admissible act, and any positive error rate, where the two-box pocket is admissible, and the ex-ante returns of the two policies differ sharply. The criterion accepts this because its standpoint for the act is the settled Aethus the deliberating agent occupies, and the ex-ante return is the return of a policy fixed at an earlier Aethus, a standpoint the criterion declines to adopt for the act. That is the normative choice the paper makes, settled evaluation over policy evaluation, and its defense beyond this statement is the burden the paper carries rather than a consequence of indexed determinacy, which selects neither standpoint. The settling rule's reply is the framework's indexing: once the full box is realized in the agent's chain it is determinate relative to the agent, and an evaluation that lets the box's contents vary with the decision-function variable treats as open what has been realized. That is the disagreement, stated as a result in one model; whether the reply persuades is the dispute between updateful and updateless evaluation, which is entered here and not settled. The divergence locates which of the functional family's verdicts the ontology underwrites and which depend on a commitment, the commitment structure of the transparent variant being pressed hardest in \cite{drescher2006good}.

\noindent\textit{The benchmarks at a glance.} Table~\ref{tab:benchmarks} collects what each benchmark requires, what the settling rule computes, and what it costs, with stakes $M = 10^{6}$ and $t = 10^{3}$ where numbers are given. Each recommendation is derived in the model from Definition~\ref{def:settle} and the stated assumptions; none is an independent intuition the rule is fitted to.

\begin{table}[ht]
\centering
\footnotesize
\setlength{\tabcolsep}{3pt}
\caption{Benchmarks under the settling rule.}
\label{tab:benchmarks}
\begin{tabular}{@{}>{\raggedright\arraybackslash}p{0.13\linewidth}>{\raggedright\arraybackslash}p{0.23\linewidth}>{\raggedright\arraybackslash}p{0.21\linewidth}>{\raggedright\arraybackslash}p{0.16\linewidth}>{\raggedright\arraybackslash}p{0.21\linewidth}@{}}
\toprule
Benchmark & Required assumptions & Settled distributions & Recommendation & Cost or limitation \\
\midrule
Opaque Newcomb & Branch-conditional accuracy $0.99$, stated as surviving deliberative settling; aggregate accuracy does not suffice & Act-dependent: full/empty $(0.99, 0.01)$ if one-boxing, $(0.01, 0.99)$ if two-boxing & One-box: $990{,}000$ against $11{,}000$ & If the prediction's marginal is held fixed instead, two-boxing wins; the verdict rests on the retained accuracy law \\
Transparent boxes & The shared model (Definition~\ref{def:shared-model}); error $\varepsilon > 0$ with every mismatch pocket positive; contents observed & Equal across acts, given the stated contents & Two-box on a full box for $\varepsilon > 0$; one-box at $\varepsilon = 0$, by admissibility & Ex-ante loss $(1 - 2\varepsilon)M - (1 - \varepsilon)t$ against policy-level one-boxing, $979{,}010$ at $\varepsilon = 0.01$ \\
Psychopath button & (J1) an act opposite to its terminal resolve is tagged; (J2) the endpoint trait law $q_1, q_0$ & Candidate-specific: psychopath with probability $q_1$ if pressing, $q_0$ if refraining & Refrain iff $q_1 > B/(B + D)$ & (J2) is an assumption about settling; ratifiability has no pure verdict when $q_0 < B/(B + D) < q_1$ \\
Smoking lesion & Resolve stated, and $C \perp L \mid R$ & Equal across acts: $P(L \mid R = r)$ & Smoke & Screening is a premise beyond the resolve's being stated \\
\bottomrule
\end{tabular}
\end{table}

\noindent\textit{Each benchmark as an instance.} Each row of Table~\ref{tab:benchmarks} is an instance of Definition~\ref{def:settle}. The opaque row has an act-to-contents coupling binding under deliberation; in the transparent row the coupling runs from the policy component $\pi_{\mathrm{full}}$ to the contents, the variable settled is the component the observation selects, and the contents enter as evidence; at $\varepsilon = 0$, when the observed opaque box is full, the two-boxing settlement has a zero normalizer and is inadmissible under the model. The psychopath button has the trait's kernel take the terminal act as input and bind under deliberation, which is what (J2) states in the operator's terms. The lesion has the lesion feed the resolve and the resolve feed the act, the resolve stated and no coupling leaving the act. The supplementary script \texttt{settling\_benchmarks.py} runs each through the definition, reproducing the table's values, shows the opaque verdict reversing when the accuracy coupling is not binding, and computes the collider example of Section~\ref{sec:invariance}.

\section{Egan's button: evaluation at the terminal resolve}
\label{sec:egan}

\noindent Egan's case challenges a pure-act ratifiability requirement, and it is run here rather than cited \cite{egan2007some}. A button, if pressed, kills every psychopath; the agent holds three beliefs --- that a world without psychopaths would be better, provided the agent is not one; that only a psychopath would press; and that they are probably not one. Evidential reasoning declines cleanly: pressing is strong evidence of being among the purged. Causal decision theory fails in two stages. Naively, the trait is causally independent of the act, so low prior credence licenses pressing --- the monstrous verdict; and with ratifiability added, deciding to press raises the credence and makes pressing look terrible given the decision, while deciding to refrain lowers it and makes pressing look good given \emph{that}. In the regime where the baseline credence lies below the stakes threshold, neither pure act is ratifiable (Theorem~\ref{thm:egan}, case (ii)); when the baseline already exceeds that threshold, pure refrain is ratifiable, and the absence of a ratifiable pure act belongs to the interior regime alone. The equilibrium accounts answer with deliberational equilibria \cite{arntzenius2008no, joyce2012regret}, dynamics of the kind already credited above \cite{skyrms1990dynamics}.

\textit{The assumption the analysis isolates.} The present run begins from the structure the case's engine has: a \emph{self-typing blank} whose coupling routes through the deliberation itself. Compare the lesion: there the mediator --- the stated taste --- screens act from trait, Reichenbach's certificate applied at the mediator. Here the mediator is the agent's own resolve --- the very quantity being settled --- and the description must be exact, because the lesion analysis above was built to keep inclinations and choices operatively separate, and that separation must survive the transfer. It does, in two clauses.

\textit{The two clauses.} The screen itself is untouched: at any \emph{fixed} resolve grade $r$, the books state $r$, and the stated resolve screens exactly as the metatickle doctrine above holds. What the case adds is jointness: an act opposite to its terminal resolve is a tagged combination, pressing showing that the terminal resolve was to press, so the act is realized only in the pockets that carry its own terminal resolve; and, by the endpoint trait law of Proposition~\ref{prop:egan-self}, a property of the settling operator rather than a consequence of the first clause, it is scored at that terminal resolve and not against the marginal credence $q(r)$ at the current one. What Egan's case changes is only the mediator's \emph{status}: the taste was exogenous, fixed while deliberation ran over it, so the screen could be applied once and the verdict read off; the resolve is endogenous, identical with the settling's own grade, so the screen holds pointwise at every $r$ while no deliberation-independent $r$ exists to condition on once and for all. Egan's case is the lesion with an exogenous mediator exchanged for an endogenous one, which is why standard ratifiability, the doctrine that conditions on the decision, and self-settled evaluation can return different verdicts here, in the interior regime. The production-mode test confirms the coupling's grade: pressed by an external randomizer, the act would raise no credence in the trait; the binding runs through the deliberative route, the tickle family's signature, and the case is prized in the literature for resisting the tickle defense \cite{egan2007some}. What is claimed here is the mechanism, the same certificate with its conditioning variable relocated into the settling. So the framework evaluates where it always evaluates: at the books the settling produces.

\begin{proposition}[The Egan verdict, on self-settled evaluation]\label{prop:egan-self}
Distinguish the \emph{current resolve} $r \in [0,1]$, the present state's weight on the press value of the act attribute, from the \emph{terminal resolve} $\tau \in \{0, 1\}$, the value a completed act settles, and let $\mathrm{Settle}(A, a)$ set the terminal resolve to $a$'s value. Assume
\begin{enumerate}[label=(J\arabic*), leftmargin=2.4em, itemsep=0.1em]
\item an act opposite to its terminal resolve is a tagged combination; and
\item the \emph{endpoint trait law}: settling carries the trait weight to the terminal resolve's, $\Pr[\text{psychopath} \mid \mathrm{Settle}(A, \text{press}),\, r] = q_1$ and $\Pr[\text{psychopath} \mid \mathrm{Settle}(A, \text{refrain}),\, r] = q_0$ at every current $r$---a property of the settling operator, not a consequence of (J1).
\end{enumerate}
With payoffs $-D$ and $B$ for pressing as a psychopath and as not, $0$ for refraining, and $q_c = B/(B+D)$:
\begin{enumerate}[label=(\roman*), leftmargin=2.4em, itemsep=0.1em]
\item \emph{Static verdict.} $V_A(\text{press}) = B - (B+D)\,q_1$ and $V_A(\text{refrain}) = 0$: refrain when $q_1 > q_c$, press when $q_1 < q_c$, either act when $q_1 = q_c$.
\item \emph{Dynamics.} The unprojected driving term $B - (B+D)\,q_1$ is constant in $r$, and the projected map $\dot r = \Pi_{T_{[0,1]}(r)}\big[B - (B+D)\,q_1\big]$ carries every resolve to $0$ when $q_1 > q_c$ and to $1$ when $q_1 < q_c$, vanishing at the boundary it reaches. When $q_1 = q_c$ the driving term vanishes, every resolve is an equilibrium, and the dynamics complete no act; the static comparison then permits either.
\item \emph{The property obtained.} The verdict is \emph{self-settled optimal}: each act is evaluated at its own settled state, and the chosen act maximizes that evaluation. This is not standard ratifiability, which evaluates the alternatives at the chosen act's settled state; the two differ exactly when an act that loses at its own settled state would win at the chosen act's.
\end{enumerate}
\end{proposition}
\begin{proof}
(i): by (J1) pressing is realized only with terminal resolve $1$, and by (J2) the trait weight there is $q_1$, so $V_A(\text{press}) = -D q_1 + B(1 - q_1)$; refraining pays $0$ in every pocket. (ii): the bracket does not depend on $r$, so its sign is constant, and the projection carries the flow to the boundary it points toward, or leaves every point fixed when the bracket is zero. (iii): the static rule compares $V_A(a)$ across acts, each computed on $\mathrm{Settle}(A, a)$, which is the definition. Standard ratifiability of the chosen act $a^{\ast}$ instead compares every alternative at the one settled state $\mathrm{Settle}(A, a^{\ast})$, and the example following Theorem~\ref{thm:egan} shows the two coming apart.
\end{proof}
\noindent\emph{Why the endpoint trait law.} The proposition's arithmetic is forced once (J1) and (J2) hold, and (J2) is where the substance lies, so its motivation is given here as a motivation and not as a derivation. The resolve is the agent's makeup as indexed by resolving to press: a makeup status-type, not a particular total makeup. In the framework's terms it is a parent Aethus, stating that the resolve is to press and leaving the rest of the agent's makeup blank, whose refinements are the specific makeups that press. The case's trait law couples the trait to that status-type, not to the momentary weight deliberation places on the act: $q_1$ is the trait's weight within the parent Aethus, its average over the specific makeups that refine it, and $q(r)$ is an average over the makeup Aethae a partial resolve still leaves open. A completed act reveals the status-type, and settling the act conditions on the revealed status-type. That alone does not give the trait weight $q_1$ whatever the prior lean: in ordinary probability $P(X \mid T = 1, r)$ can depend on $r$ although $T = 1$ is a legitimate coarse-grained type. (J2) is the conditional-independence premise that, once the terminal resolve is specified, the current resolve and any retained information carry no further information about the trait; a model in which they do is a different model, and the paper adopts (J2) as a stated premise. The contribution is accordingly a proposed endpoint evaluation policy with an indexed interpretation, and not a proof that rival evaluations combine impossible objects: (J1) prohibits an actual completed act from disagreeing with its terminal resolve and does not prohibit representing a hypothetical alternative while deliberating, so a rival calculation is not ``tagged'' but owed an argument against its counterfactual policy. A reader who takes the vignette's coupling to attach to momentary deliberative weight rejects (J2), and with it the verdict.

\noindent\emph{The relata of the feedback account.} The challenge to the feedback account concerns its relata: does the trait probability belong to the agent's current deliberative state, or to the status-type obtained by settling the candidate act? The rule evaluates the latter. The feedback model is not thereby undefined; a rival may defend its pairing of the current state's trait distribution with a completed act as its evaluation policy. And (J2) is a commitment in its own right, entered as an assumption rather than folded into well-definedness, since identifying pressing with the terminal resolve does not alone fix the trait distribution at $q_1$. The contribution claimed is that the indexing makes this dependency explicit, not that other frameworks cannot formulate the question.

\noindent\emph{What is new, against the nearest positions.} The status-type reading has relatives: Eells's tickle defense \cite{eells1982rational}, Jeffrey's ratifiability \cite{jeffrey1983logic}, the deliberational-dynamics tradition \cite{skyrms1990dynamics, arntzenius2008no, joyce2012regret}, Price's agent's perspective \cite{price2012causation}, and functional decision theory's question of which algorithm the agent is \cite{soares2017fdt}, the nearest kin to asking which status-type. Ratifiability and the dynamics tradition condition on, or track credence in, a decision event; once the trait is coupled to how strongly the agent leans at a moment, the target moves as deliberation moves, and the equilibrium accounts respond to that instability. Joyce's response is a substantive defense of causal decision theory rather than a repair: the agent must process the information deliberation itself generates before any choice is binding \cite{joyce2012regret}, and the present proposal differs from it in where the trait law is read, at the terminal resolve rather than along the deliberation.

\noindent The proposition is the paper's own theorem about the case. Appendix~\ref{sec:comeapart} compares it with an illustrative comparison dynamics, the projected additive flow $\dot r = B - (B+D)\,q(r)$ with endpoint projection, which scores pressing against the trait credence at the current resolve; the theorem there is about that specified flow and not about every deliberational-dynamics account, a logistic variant $\dot r = r(1-r)[B - (B+D)q(r)]$ having both endpoints as fixed points even where the interior equilibrium attracts every interior start. The affine coupling $q(r) = q_0 + r(q_1 - q_0)$ is what total probability gives when $r$ is the current weight on the terminal resolve and (J2) is read as conditional sufficiency within one joint model, $P(\text{psychopath} \mid r) = r q_1 + (1 - r) q_0$; the theorem's general strictly monotone $q$ covers a broader comparison family than that model.

\noindent\emph{The vignette, and the books.} The verdict is indexed to the vignette's stipulated attribute stock: a real deliberator's books carry entries the idealization fences out, above all the option value of an irreversible act over blanks whose couplings are unknown, and the verdict is owed recalibration wherever the stock grows.

\noindent\emph{Dynamic deliberative settlement, as an illustration.} The static rule of Principle~\ref{prin:rule} governs; a settling map describes one way its selections are reached, and it is supplied for this case and not for the whole model class. For a resolve state $r$, write $\dot r = \Phi_A(r)$ with equilibrium set $\mathrm{Eq}(A) := \{r : \Phi_A(r) = 0\}$, pure acts being the boundary equilibria $r \in \{0, 1\}$. Under (J1) and (J2) the map moves the resolve toward the act the static comparison favors and its equilibria are the boundary points that comparison selects, so the theory's output is an act, and the rule is a complete act-selection rule; at exact indifference the field vanishes, every resolve is stationary, and the dynamics complete no act, so the correspondence between the static rule and the dynamics is a correspondence of verdicts, not a claim that deliberation terminates. A map that scores an act against the trait credence at a partial resolve lets act and resolve come apart; its interior equilibria, where they arise, describe deliberation on that reading only, which Appendix~\ref{sec:comeapart} sets beside the paper's own.

\section{Why evaluate each act at its own settled state, and how the rule compares}\label{sec:defense}

\noindent The rule needs a defense of its own, because the ontology motivates it without selecting it, and one question carries the defense: \emph{why should an agent adopt this rule when a competing policy performs better ex ante in the transparent case?} The answer has two parts, one about which couplings bind and one about which information state an act is evaluated at, and the rule stands or falls with them.

\textit{Why the couplings bind.} The binding assignments are part of the argument, not implementation details: the verdicts of Table~\ref{tab:benchmarks} follow once they are supplied, so they are defended case by case. The model stipulates that the branch-conditional accuracy law remains binding under deliberative settling, and the stipulation is defended as the paper's reading of the problem rather than proved. The predictor's accuracy does not by itself settle it: accuracy conditional on the agent's actual deliberative choices does not determine which dependencies survive a counterfactual evaluation, and that is part of what the evaluation policies dispute. The reading takes the predictor's tracking of the deliberating agent to hold of the agent's own settling of the act; a reader who severs the coupling instead, as the causal rule's surgery does, evaluates a different model of the problem, on which that rule's two-boxing verdict follows. The assignments are therefore defended reading by reading and are not supplied by the ontology; where a vignette leaves the reading open, the rule's verdict is correspondingly open. Settling uses the agent's own indexed structure for a reason the causal companion makes precise: dependencies established at other indices neither transfer into it nor compose with one another without a common admissible index and a composition rule \cite{AethusCausality}.

\textit{Why the transparent-box cost is worth paying.} The competing policy's ex-ante advantage belongs to a policy settled before the observation, at an information state in which the box's contents are open, and this rule secures that advantage at that state: settled there, it recommends the one-boxing policy (below). This rule's requirement is its own, and it is stated as a constraint on evaluation: \emph{a candidate settlement must preserve the information already stated at the agent's current index}. Under it the full box is stated and the advantage is not available. The constraint is not proved by the definition of the rule, and saying that an evaluation which varies the box's contents concerns an act the agent is not performing does not justify it, since a rival uses exactly such a counterfactual comparison to decide what to perform. The argument for adopting the constraint has four parts. First, the conception of act evaluation it serves: an act is evaluated by what its performance would make true of the agent's own situation at the index of choice, the law of the state its settling induces, so that the evaluation never assigns an act a law the agent could not be in after performing it; this is evaluation at the agent's evidence, with the evidence preserved. Second, why that conception is given priority: the stated contents are the agent's evidence about the situation they are in, and an evaluation that lets them vary answers a question about a different situation, the one an agent at the earlier index was in, which is the question policy evaluation answers; the rule holds that the act at the box is the object to be evaluated, and that its worth is fixed by the situation its performance induces rather than by what a policy fixed earlier would have produced on average. Third, why the acknowledged costs are acceptable: the preference reversal between the earlier and the later index is the standard updateful--updateless tradeoff, and the rule answers it with commitment rather than with a demand for dynamic consistency, an agent who values the ex-ante return being told to commit before looking, the commitment being worth what the model says it is worth. Fourth, what the argument leaves available: policy-level, updateless evaluation is a different conception, answering which policy fixed before the observation does best, and nothing here refutes it; the disagreement is about which question act evaluation should answer, and the rule's case is that the question about the act at the agent's own index is the one a theory of acts is for. A sequential protocol makes the commitment concrete. Before the prediction and the observation, the agent can purchase an enforceable commitment to a policy; the predictor tracks the resulting policy under the stipulated error law; after the observation, an uncommitted agent applies the act rule. At the canonical numbers the one-box policy is worth $990{,}010$ and the two-box policy $11{,}000$, so the commitment is worth $979{,}010$ at the earlier state, and the protocol states the preference reversal plainly: the earlier state favors the one-boxing policy, and the later, full-box state favors two-boxing whenever both acts remain admissible. Enforceable commitment is acceptable on the account's own terms because the committed agent's later act is no longer a candidate evaluated at the full box but a stated consequence of the earlier settlement.

The defense then has three parts: what the rule preserves, what it gives up, and where it parts from its strongest rivals.

\textit{What it preserves.} The policy adopts one requirement: each candidate is evaluated under the law of the information state its own performance would settle, so that no act is assigned a law the agent could not be in after performing it. That is a requirement the policy adopts, not a demonstration that rival evaluative counterfactuals are incoherent; the rivals adopt other requirements, compared in the uniform comparison below. Dominance is not abandoned: wherever the settled laws coincide across acts, which Lemma~\ref{lem:invariance} characterizes, pointwise dominance transfers to expected utility (Corollary~\ref{cor:dominance-transfer}), and the rule then agrees with every dominance argument. And observed content is held fixed: what the agent has read is stated, and no evaluation reopens it.

\textit{What it gives up.} Two things. First, the fixed-distribution dominance argument in exactly the cases where the settled laws depend on the act. Blankness alone does not defeat dominance: a prize held at one probability $q$ under both acts is blank and leaves two-boxing ahead by $t$. What defeats the argument's applicability is an act-dependent settled distribution, which arises only when a stated coupling binds across the act's settling; the rule then compares whole residues instead. Second, ex-ante payoff in the transparent case. The rule two-boxes on a full box at every positive error, and an agent known to do so is less often shown a full box: against policy-level one-boxing the ex-ante loss is $(1 - 2\varepsilon)M - (1 - \varepsilon)t$, which is $979{,}010$ at $\varepsilon = 0.01$. The rule accepts that cost for the reasons just given. An agent able to commit before the observation is at a different state, the one in which the policy's effect on the prediction is open; settling the policy there, with the accuracy law binding on it, gives the policy values of Proposition~\ref{prop:shared-model}, so the same rule recommends the one-boxing policy whenever $\varepsilon < (M - t)/(2M - t)$. What the rule declines is to evaluate the act at the box at any state but the one the agent occupies there; policy-level evaluation answers a different question, about the policy at the earlier state.

\paragraph{The comparisons, uniformly.}\label{sec:locating} The rule's relation to the existing positions is set out with the model and the information held fixed, since the benchmark verdicts alone do not separate it from its nearest neighbors: an evidential evaluator supplied with the same screening and conditional-probability assumptions returns the four verdicts of Table~\ref{tab:benchmarks}, and the updateful variant of functional decision theory characterized by MacDermott, Everitt and Belardinelli one-boxes in opaque Newcomb and two-boxes in transparent Newcomb, the pair of verdicts this paper reaches \cite{macdermott2023characterising}. Each comparison below names the example, if any, on which the rule and the comparator part.

\textit{Against causal decision theory and the marginal-preserving rule.} The disagreement concerns what is held fixed when candidate acts are evaluated: causal decision theory's dominance argument requires that Box B's contents be a fixed fact across which the agent's act cannot range, and tense-free indexing motivates the settling policy that lets them range without by itself selecting it. Under the adopted branch-conditional settling policy the prize distribution differs between the candidate acts, so the usual fixed-distribution dominance argument does not apply: there is no single fixed state of Box B, held constant across acts, over which taking Box A yields a guaranteed extra thousand; blankness alone would not defeat dominance, as the conclusion's fixed-$q$ case shows. The causal decision theorist and the present proposal disagree about the settling rule, the marginal-preserving against the branch-conditional; whether the past is determinate for the deciding agent is what makes the branch-conditional rule available and motivates it, not what selects it. The separating example is opaque Newcomb with the accuracy coupling stated as binding: the settled prize laws differ across acts, and the rule one-boxes where the marginal-preserving rule two-boxes. The marginal-preserving rule is coherent; it treats the evaluation of an act as an intervention that leaves the prediction's generating mechanism untouched, and the disagreement turns on one modeling commitment, whether the predictor's branch-conditional accuracy survives the deliberative mode of the act's production: where it does, the settled laws differ and the rule one-boxes, and where it does not, they coincide and the rule two-boxes with the causal theorist.

\textit{Against evidential decision theory.} The relation is convergence-with-difference. Both reach one-boxing, and both regard the two-boxer's ``the prediction is already fixed'' as the error. But the diagnosis differs: evidential decision theory grants that the prediction is fixed and locates the two-boxer's mistake in ignoring the act's evidential bearing on that fixed fact, whereas the present proposal locates it one level earlier, in treating the prediction as fixed at all. The verdict is shared; the reason is not, and the operational difference is the detachment of nonbinding couplings. Two examples separate the rules on the same record. A predictor whose accuracy is stated as robust to habit alone, $m(e) = \{\text{habit}\}$: an agent whose act issues from habit one-boxes under the rule, $990{,}000$ against $11{,}000$ at $\varepsilon = 0.01$, while a deliberating agent two-boxes, $500{,}000$ against $501{,}000$, the coupling being detached under deliberation, whereas the evidential rule, conditioning on the act in the record, one-boxes for both. And the unscreened lesion: with the desire blank and not stated, the fix severs the act's dependence on the desire, the settled lesion law is the prior under both acts and the rule smokes, while conditioning on the act in the observational law refrains; the supplementary script computes both.

\textit{Against the functional family.} The closest convergence is with \emph{timeless} and \emph{functional} decision theory (Yudkowsky \cite{yudkowsky2010tdt}; Yudkowsky and Soares \cite{soares2017fdt}; Levinstein and Soares \cite{levinstein2020cheating}). Those theories reach one-boxing by having the agent reason as though their decision procedure determines the prediction, through the predictor's model of that procedure, a simulation, notes or a mental model---so that choosing one-boxing ``logically'' brings about the million by fixing what the predictor's model of the agent outputs. This is structurally the right shape, and the present proposal agrees with its verdict and much of its spirit. Against the updateful variant no benchmark separates the two, and none is claimed: once a binding coupling is identified with a subjunctive dependence of the decision procedure, the verdicts coincide on every case run here, and the operator differs in form, indexing binding to the mode of the act's production and detaching nonbinding couplings jointly by act-averaging rather than by removing the node's influence, with no separating case exhibited. Against the updateless, policy-level comparator the separating case is transparent Newcomb on a full box, Proposition~\ref{prop:shared-model}, where the disagreement turns on what the evaluative counterfactual holds fixed, the observed contents or the error law with the prediction varying; agreement on opaque one-boxing is therefore not equivalence, and the transparent case is where the rule parts from that comparator and from nothing else in the family. What those theories leave unanchored is the ontological status of the ``as though'': in what sense does a present decision determine a past placement, if not causally and not merely evidentially? The present proposal supplies that anchor. The decision determines the prediction neither by an as-if stipulation nor by backward causation, but because determinacy is observer-indexed and the deciding agent, in settling the choice, settles which of the consistent states---each with its own already-consistent past---is the one they occupy. The ``logical'' determination that timeless decision theory posits is, on this reading, the deductive closure of an information state doing visible work. The contribution against the functional family is accordingly of the second kind: not a new verdict function but an interpretation and justification for a procedure within an existing family, together with the mode-indexed binding structure as an explicit input and the representation of Proposition~\ref{prop:aux}; the contribution against the evidential and causal rules is the formal one the separating cases show.

\paragraph{What the operator does differently, with the model and information held fixed.} On the model class of Definition~\ref{def:settle} the rule is a hybrid of conditioning and intervention: the fix step severs the act's dependence on its parents, as $\mathrm{do}(C = a)$ does; a coupling that binds in the evaluation mode is kept reading the settled act, as conditioning keeps it; a coupling that does not bind is replaced by its act-average under the original production entry, as a stochastic intervention on the edge from the act to that output does. Where every act-relevant coupling binds, the settled laws are the act-conditional laws after the fix, which are the evidential laws wherever the act's parents are stated or absent; where none binds, they are the laws of the marginal-preserving rule. The differences from each rival are produced by the binding assignments and not by the arithmetic. Standard ratifiability evaluates every alternative at the chosen act's settled state; this rule evaluates each at its own, and in the psychopath-button case, when $q_0 < B/(B + D) < q_1$, ratifiability finds no ratifiable pure act while the rule refrains, a difference of evaluative commitment, not of computation.

\paragraph{The Stated-Relata Principle.} The lesion verdict consumed one discipline, the \emph{Stated-Relata Principle}: the framework's requirement that conditioning be on well-defined attributes, applied to couplings. Three things are distinct: an attribute's being available in the state's vocabulary, its value's being stated, and a relation involving it being stated. A coupling can be stated while the values of its relata are blank, as the accuracy coupling is in opaque Newcomb with the act and the prediction both unresolved; what the principle requires is that a coupling's relata be attributes the state carries. Three questions are then kept apart: whether a proposed relatum has a defined representation in the state, whether an observational dependence holds between the represented attributes, and whether that dependence survives the specified settling. The principle bears on each, and a failure at the third does not establish a failure at the first.

\textit{The lesion as exhibit.} The lesion is the principle's cleanest exhibit once one further distinction is drawn. ``I smoked'' is a perfectly formed stated attribute; nothing bars the act from the books. What the classical argument needs is a different object --- ``I smoked, \emph{in the sense that carries the lesion's evidential charge}'': the act as psychological symptom --- and the association that object carries, between the act and the lesion, is representable in a probabilistic model, the rival's variable being well defined; what the argument needs is that the association survive the specified settling, and it does not, the coupling's binding relatum being the inclination, an attribute \emph{about} the mental layer rather than a dependency on it. The argument thus cites an association that a probabilistic model represents but that does not survive the specified settlement operation, and the rule denies it authority over the act's evaluation on that ground alone: the rival's variable is well defined, and the disagreement is about the operation, not about an ontological defect in the represented variable. Production-mode robustness, the property of binding under every mode of act-production, is a diagnostic for the operation and not a second principle. Perturbing how an act is produced varies which couplings bind while holding the stated entries fixed, so a dependence that does not survive the perturbation is one the specified settling does not retain under the perturbed mode; the diagnostic is stronger than any verdict needs, since placement under the stated mode of act-production already selects the couplings that bind, and a coupling binding under deliberation is retained for a deliberating agent whether or not it would survive a randomizer.

    \textit{Is this causal dependence in new vocabulary?} It is the objection a decision theorist will reach for first, and it deserves a direct answer with an exhibit rather than a denial. Consider two faithful implementations of the same function on the same inputs---the same algorithm run on silicon here and worked out with pen and paper there. Their outputs are perfectly correlated; the correlation is robust to changes of implementation, surviving any intervention on substrate, timing, or location, which is a weaker property than robustness under replacement of the decision-producing process by a randomizer and is all the exhibit claims; and there is no causal arrow in either direction, neither instantiation being upstream of the other and no common cause screening the correlation off, since the dependence is a fact about the procedure rather than about any channel between the two runs. Act-robustness registers this coupling as it stands. A causal modeller may reply that the two runs share an exogenous input---the procedure's specification---and call that the common source; the exhibit is not offered to defeat that reply, and the paper does not need it to. The claim the exhibit supports is the weaker and sufficient one: act-robustness is not \emph{directed act-to-outcome causal influence}, and need not be represented by such an arrow, since the coupling here has no direction and no channel between the runs. Whether every act-robust coupling admits \emph{some} causal representation, common-source or otherwise, is a question the benchmark verdicts do not turn on, since those verdicts are largely coextensive with functional dependence in any case; what they turn on is that the coupling is not an arrow from the act.

    Two concessions keep the claim exact. Over the cases this paper treats, act-robustness and the functional or subjunctive dependence of the timeless theories sort Newcomb and the lesion alike, agree with the updateful variant on the transparent case, and part from the updateless comparator there, where policy-level evaluation one-boxes and act-robustness, the prize already stated, two-boxes at positive error (Proposition~\ref{prop:shared-model}). What differs is the grounding. Functional dependence is read off a shared computation, which the functional presentation models by the predictor's possession of a model of the agent's procedure; act-robustness is read off the stated content of an indexed information state and asks only whether the correlation's relatum is an attribute on the books, a demand that applies to a predictor whose method is unspecified. A reader who holds that all such couplings reduce to functional dependence may take the criterion as an operational test for it and lose nothing this paper uses. At its defensible strength, act-robustness is an implementation-independent diagnostic for the dependence the functional theories posit, with a different grounding and a wider domain, and not a dependence relation extensionally distinct from it on the benchmark family.

\paragraph{Betting on the past.} Ahmed's case is the tradition's cleanest attempt to make the fixed past itself the lever \cite{ahmed2014evidence}. Stipulate determinism, and let $P$ say that the world's state at some pre-natal time, together with the laws, entails the agent's now taking Bet~1; the bets are arranged so that taking Bet~1 makes $P$ true and taking Bet~2 makes it false, with the payoffs set so that the causalist --- holding the past fixed and screening the act from $P$ --- mishandles a bet the agent literally cannot lose. The present framework's verdict is immediate, and it is reached with no new machinery: $P$'s subject is a past attribute \emph{blank on the agent's books} --- the microstate is undeduced --- and its entry is coupled to the act at entailment grade, the strongest binding mode there is, robust to every mode of production by the stipulated laws themselves. So the perturbation of the act restates $P$ across the entire surviving residue, the case runs as ordinary act-dependence, and the guaranteed win is taken. What this adds to the evidentialist's matching verdict is a diagnosis and a prediction. The diagnosis: the causalist's error is the shared-ledger construct once more --- a determinate, act-independent value assigned to $P$ \emph{on the agent's own books}, exactly the fixed-past premise this paper deletes, with tense-free blankness making the pastness of $P$ do no work at all. The prediction, checkable against the literature's own reception: the case should be structurally identical to Newcomb with the predictor replaced by the laws --- one coupling, one criterion, one verdict --- and it is. That the tradition experienced the past-tense version as a \emph{separate} embarrassment is itself evidence that the operative variable was never the tense.

\paragraph{The spectator objection.} One further pressure point of Nozick's appears to favor the causal camp decisively: the friend behind the boxes \cite{nozick1969newcomb}. A friend has already seen Box B's contents before the agent decides; from the friend's standpoint the contents are fixed, and the friend, wanting the agent to do well, urges taking both boxes, since whatever is in Box B the extra thousand can only help. On the present proposal the friend's reasoning is correct relative to the friend's state, in which the contents are stated, and the error is in importing it into the agent's decision: relative to the agent's state the friend's observation is itself blank, the friend occupying a superposition over having-seen-the-million and having-seen-nothing, correlated with the still-open prediction, the structure that arises whenever one observer's measurement is unread by another \cite{Aethus}. One sentence makes the diagnosis formal: the friend's urging is content-invariant by construction, counselled whatever was seen, so its arrival prunes nothing in the agent's state; the friend's reasoning is valid in the friend's state and the friend's utterance uninformative in the agent's. A friend who instead says \emph{I looked, and it is empty} makes a content-bearing utterance that enters the agent's state and settles the contents there, which is the third grade's account of how determinacy arrives. The case illustrates the ontology's commitment, that what is determinate for an agent is fixed by what the agent's state states, with no further fact behind the states to adjudicate between them, and it does not establish that objective determinacy merges information illicitly: under a fixed-world probability model the content-invariant urging transmits no information about the contents either, so the classical information-theoretic explanation covers the friend's advice as well. The friend's two-boxing verdict is the positive-error one; at the perfect endpoint an agent who saw what the friend sees falls under the admissibility rule of Corollary~\ref{cor:transparent-sort}, and the friend's content-invariant counsel would be inadmissible whenever the million is visible.

\section{Conclusion}
\label{sec:conclusion}

\noindent Newcomb's problem has been durable because both of its standard resolutions accept that a past prediction is a present determinate fact. This paper has proposed dropping that premise. On observer-indexed determinacy an unread past outcome is a blank entry of the deciding agent's information state, unsettled regardless of its position in time, and a specified settling rule, evaluating each act at the state its own settling induces, returns the benchmark verdicts of Table~\ref{tab:benchmarks}: one-boxing without news value or backward causation, converging with the functional family in the opaque case, agreeing with its updateful variant in the transparent case, and parting there, at a stated cost, from the updateless, policy-level comparator.

None of this is offered as a proof that one-boxing is correct. The two-boxer who keeps tensed determinacy is not refuted; the indexing motivates the evaluation objects without selecting the policy, since a prize held at one probability $q$ under both acts leaves the blank in place and two-boxing ahead by $t$, and the verdict rests on the settling rule together with the retained branch-conditional accuracy law. The sealed box is, moreover, a record-rich case: on the wider framework's own criterion the agent's blank is an exclusive superposition, for which ``genuinely open relative to the agent'' and ``determinate but unread'' issue identical statistics \cite{AethusBrief}, so what the proposal offers here is architectural, one treatment of determinacy applied uniformly, and not an empirical adjudication.

\noindent\textit{What is established.} On finite acyclic graphical models with supplied roles and binding modes, the settling rule is well defined (Proposition~\ref{prop:settle-exists}) and returns the verdicts of Table~\ref{tab:benchmarks} under the assumptions listed there. The ontology removes the tensed premise; the policy is adopted, and Section~\ref{sec:defense} defends it with its costs, the transparent-box loss among them. The d-separation condition of Corollary~\ref{cor:dsep} is the graphical sufficient condition for across-act invariance; whether a broader graphical criterion characterizes it, and how the rule extends beyond acyclic models, remain open, and the probability calculation of the settled laws is the authoritative test.

\section*{Provenance} The paper was drafted and revised under adversarial review directed by the author, whose rulings fixed every substantive choice; where review supplied a formal device rather than a translation, the session record names it. % R823 (Fable 5.1, Oct 18 2026): the dates aligned with the long paper's record (reached by mid-2020; first written June 9, 2020; separated into Nexic reasoning in 2021). Original: dates to 2019--2020, where it was first stated in a question of natural history, whether the extinction of the dinosaurs was a prerequisite for observers like us,
The typicality principle of Appendix~\ref{sec:anthropic} was reached by mid-2020, in a question of natural history, whether the extinction of the dinosaurs was a prerequisite for observers like us, its first written record a script of June 9, 2020; Nexic reasoning as a formalism dates from 2021, when the principle was separated from that question and stated generally; and the Aethic and Nexic manuscripts this paper draws on were complete before AI-assisted drafting began in June 2026 \cite{Nexus}. The apparatus was generalized from the Doomsday argument \cite{gott1993implications, leslie1996end} independently of the anthropic-principle literature; the documentary chain is in the companion \cite{Nexus}, and the natural-historical application is developed in \cite{NexicFoundation}.

\appendix
\section{Why the block is detached jointly and from the act's inputs}\label{sec:withdrawn-operators}

\noindent Step (3) of Definition~\ref{def:settle} has two features, and each is forced by a minimal example.

\emph{Jointly.} Let the act $C$ be a fair bit, produced without inputs, with two non-binding couplings $Y_1 := C$ and $Y_2 := C$, and let $Y_1 = 1$ be stated. The joint kernel $K(y_1, y_2) = \sum_c \tfrac12 [y_1 = c][y_2 = c]$ puts weight $\tfrac12$ on $(0,0)$ and on $(1,1)$, so conditioning on $Y_1 = 1$ gives $Y_2 = 1$ with certainty, whichever act is settled. Detaching the couplings one at a time replaces each by the act-marginal, $Y_1$ and $Y_2$ independent fair bits, and conditioning on $Y_1 = 1$ leaves $Y_2$ fair. Separate detachment discards the dependence the shared act induced among the block's outputs; joint detachment keeps it while removing their dependence on the settled value.

\emph{From the act's inputs.} Let $W$ be a fair bit, the act produced as $C := W$, with one non-binding coupling $Y := C$, and let $W = 1$ be stated. The kernel from the act's production inputs, $K(y \mid w) = \sum_c P(c \mid w)[y = c]$, gives $Y = 1$; detaching under the act's unconditional law, $K(y) = \sum_c P(c)[y = c]$, makes $Y$ fair despite the stated input. Detachment from the inputs keeps the block's dependence on the act's inputs across the block's boundary.

\emph{What is not claimed.} The rule asserts across-act invariance of the settled law only under the factorization hypothesis of Lemma~\ref{lem:invariance}, verified per benchmark; the opaque case's settled distributions differ across acts by design. Nor does it preserve standing marginals: the second example shows a non-binding output's law changing with evidence that enters through an input of the kernel, and a binding coupling stating the same output can change it likewise.

\section{The comparison dynamics: scoring at the current resolve}
\label{sec:comeapart}
\noindent\textit{The current-resolve map, and ratifiability.} The comparison dynamics scores pressing against the trait credence at the current resolve, and lets act and resolve come apart; its interior fixed point, which recovers ratifiability, is derived here with its conditions, since the conditions matter. Fix the payoffs: pressing while a psychopath costs $D > 0$ (the agent dies), pressing while not a psychopath yields $B > 0$ (the world improved, the agent alive), and refraining yields $0$ in either case. Let $r \in [0,1]$ be the resolve toward pressing, the weight the agent's state places on the press value of the act attribute, and let resolve-luminosity give the coupling as a function $q(r) = \Pr[\text{psychopath} \mid \text{resolve } r]$, continuous and strictly increasing, with $q(0) = q_0$ the baseline and $q(1) = q_1$ the credence at full resolve. Then
\begin{equation}
\label{eq:egan-eu}
\mathrm{EU}(\text{press} \mid r) \;=\; -D\,q(r) + B\,(1 - q(r)) \;=\; B - (B + D)\,q(r), \qquad \mathrm{EU}(\text{refrain} \mid r) \;=\; 0 .
\end{equation}
Two maps present themselves. The deliberational dynamics of Skyrms \cite{skyrms1990dynamics} written on the act attribute scores pressing against the marginal credence $q(r)$ at the current books,
\begin{equation}
\label{eq:egan-map}
\dot{r} \;=\; \mathrm{EU}(\text{press} \mid r) - \mathrm{EU}(\text{refrain} \mid r) \;=\; B - (B + D)\,q(r),
\end{equation}
clipped to $[0,1]$; a fixed point is a resolve at which the books, stating $r$, give the agent no reason to move it---ratifiability, stated on the settling. That map evaluates pressing by the trait distribution at the current resolve, whereas the paper's policy evaluates pressing by the distribution at its own terminal resolve; jointness constrains the realized act--resolve combinations, and the endpoint trait law, (J2), supplies the further commitment that distinguishes the two evaluations. On the paper's reading, by the endpoint trait law, pressing's worth is $B - (B+D)\,q_1$ at every $r$ and refraining's is $0$, so the unprojected driving term is $B - (B+D)\,q_1$, constant in $r$: the resolve settles at the boundary the static comparison selects, refrain when $q_1 > q_c$ and press when $q_1 < q_c$, and at exact indifference, $q_1 = q_c$, the field vanishes, every resolve is an equilibrium, and either act is permitted (Proposition~\ref{prop:egan-self}). The theorem below works that map as the comparison dynamics; its regimes are what the case's literature turns on, and the paper's own verdict is the jointness one.
\begin{theorem}[The Egan fixed point, on the current-resolve map]
\label{thm:egan}
Let $q_c = B/(B+D)$. The settling map \eqref{eq:egan-map} has exactly one fixed point $r^{\ast}$ in $[0,1]$, and it is globally attracting:
\begin{enumerate}[label=(\roman*), leftmargin=2.4em, itemsep=0.1em]
\item if $q_0 \ge q_c$, then $r^{\ast} = 0$: pure refrain is ratifiable;
\item if $q_0 < q_c < q_1$, then $r^{\ast} = q^{-1}(q_c) \in (0,1)$: the fixed point is interior, and neither pure act is ratifiable;
\item if $q_1 \le q_c$, then $r^{\ast} = 1$: pure press is ratifiable.
\end{enumerate}
In case (ii) with the affine coupling $q(r) = q_0 + r(q_1 - q_0)$, $r^{\ast} = (q_c - q_0)/(q_1 - q_0)$. As $D/B \to \infty$, $q_c$ falls, $r^{\ast}$ decreases through case (ii), and once $q_c \le q_0$ the equilibrium is case (i)'s boundary $r^{\ast} = 0$; so $r^{\ast} \to 0$ across regimes, the interior formula not being followed past $q_c = q_0$.
\end{theorem}
\begin{proof}
The right-hand side of \eqref{eq:egan-map} is $B - (B+D)q(r)$, strictly decreasing in $r$ since $q$ is, and it vanishes exactly where $q(r) = q_c$. If $q_0 \ge q_c$ it is non-positive on all of $[0,1]$, so $r$ falls to $0$ and stays; if $q_1 \le q_c$ it is non-negative throughout, so $r$ rises to $1$ and stays; and if $q_0 < q_c < q_1$, continuity and strict monotonicity of $q$ give a unique $r^{\ast} = q^{-1}(q_c)$ in $(0,1)$ at which it vanishes, positive below and negative above, so the dynamics converge to $r^{\ast}$ from any start. The global attraction has a one-line Lyapunov proof: with the boundary convention that the flow is projected onto $[0,1]$---$\dot r = \Pi_{T_{[0,1]}(r)}[B - (B+D)q(r)]$, the vector field clipped at the endpoints---take $V(r) = \tfrac12 (r - r^{\ast})^{2}$; then $\dot V = (r - r^{\ast})\big(B - (B+D)q(r)\big) < 0$ for $r \ne r^{\ast}$ in the interior, since the bracket is positive exactly below $r^{\ast}$ and negative exactly above it by the strict monotonicity of $q$, and the projection at the boundary cannot increase $V$. Uniqueness in every case follows from strict monotonicity. For the affine coupling, solving $q_0 + r(q_1 - q_0) = q_c$ gives the closed form, and $q_c \to 0$ as $D/B \to \infty$ forces $r^{\ast} \to 0$ in case (ii) and case (i) to obtain once $q_c$ falls below the baseline.
\end{proof}
\noindent Two things the theorem adds. First, the interior fixed point is conditional: it exists exactly when the baseline credence sits below the stakes threshold $q_c$, so at Egan's stakes, $D \gg B$ and $q_c$ near zero, a baseline at or above $q_c$ puts the case in regime (i), where pure refrain is itself the fixed point; $D/B = 1000$ with a baseline of $10^{-6}$ still lies in regime (ii). The claim that neither pure act is ratifiable holds only in regime (ii), and there the interior point is an indifference point at which the dynamics select no pure act. Second, the fixed point is unique and globally attracting, so the resolve the couplings support is one thing. Call a resolve \emph{resolve-ratifiable} when the expected-utility gradient computed at it supplies no reason to move; a pure act is resolve-ratifiable iff it is ratifiable in the standard sense, and the interior fixed point is resolve-ratifiable and no more. This paper's verdict is not that equilibrium but the self-settled optimality of Proposition~\ref{prop:egan-self}, and the two can differ: with $B = D = 1$, $q_0 = 0.1$ and $q_1 = 0.9$, refrain is the verdict, pressing at its own settled state being worth $1 - 2q_1 = -0.8$, yet pressing scored at refrain's settled state is worth $1 - 2q_0 = 0.8$, so refrain is not ratifiable in the standard sense. The standard test makes a comparison the adopted policy does not make; (J1) constrains completed act--resolve combinations and does not rule that comparison out, so the difference is one of evaluative policy, not a proof that the comparison is illegitimate.
\begin{center}\emph{Recommendation: }\textsc{refrain}\emph{ whenever $q_1 > q_c$, press whenever $q_1 < q_c$, and either act at $q_1 = q_c$, by self-settled evaluation; the interior fixed point of the current-resolve map describes the equilibrium accounts' deliberation and is not a verdict.}\end{center}
The causalist's original verdict presses with the credences of an agent not yet resolved to press, an evaluation at one index consumed at another; against the equilibrium accounts, the differentia is that on indexed books an act is realized only with its own resolve, which with (J2) gives the case its pure verdict.

\section{A typicality restatement of the same verdict}
\label{sec:anthropic}
\noindent There is a second route to the opaque verdict, developed in the companion framework \cite{Nexus, NexicFoundation}, and it is recorded here as an interpretation rather than as an independent derivation. The framework's Accordance Principle extends the mediocrity principle, that one's own observations are typical of observers in one's reference class, from locations in space and time to the space of possible histories. Applied to Newcomb, the predictor's record is read as population-level information about the reference class the agent belongs to, and settling the act places the agent in the sub-class that performs it; on the branch-conditional accuracy retained under settling, the members who one-box collect and the members who two-box do not, and the counsel is to act as the collecting members do. That is the main text's calculation in another register: the placement supplies the same conditional payoff laws, and a reader who rejects the principle keeps the decision rule as the stated policy of Section~\ref{sec:proposal} and loses this grounding for it. The reference-class problem applies to the route as to every typicality argument, and its answer is the one the main text gives, that the sub-class is fixed by the act's settling under the stated mode of production; the principle's own statement, its distance from the anthropic literature, its finite probabilistic form, and the self-location cases it bears on are the companion's business and are not argued here.

\section{The randomization objection}
\label{sec:randomization}

The strongest objection to one-boxing is that the agent can refuse to decide at all, delegating the choice to a genuinely random device---a quantum coin, say---so that there is no decision procedure for the predictor to have modeled, and any predictive mechanism is apparently nullified. Nozick anticipated this and handled it by stipulation: an agent predicted to randomize finds Box B empty \cite{nozick1969newcomb}. As a stipulation it settles the case by fiat rather than explaining it, and it leaves open what a predictor faced with genuine quantum randomness is actually supposed to do. The framework replaces the stipulation with a mechanism that is not local to decision theory, the joint-consistency structure of Principle~\ref{prin:joint_invalidity}, of which this appendix's verdict is the decision-native instance \cite{AethusBrief}. One further assumption is in play: a predictor accurate against a genuinely quantum randomizer is, on the framework's own physics, impossible in-world --- the accuracy stipulation analyzes a fantastical premise, as Newcomb's original predictor does, and the analysis is of the vignette's internal logic, never a physical proposal.

The present framework makes an answer available, and it imports one additional stipulation beyond the bare decision rule of Principle~\ref{prin:rule}: that the predictor's near-certainty holds \emph{even against} a randomizing agent. What is derived is the consequence of that strengthened premise. (Definition~\ref{def:binding} promotes this assumption to a general definition of binding modes, on which the sorting of the cases turns.) Granting that---as the problem's own premise of a near-perfect predictor invites---the framework represents the strengthened premise as a joint-consistency constraint; it does not derive the premise, which remains the stipulation the abstract names.

\textit{The joint system.} Consider the joint system of three attributes: the coin's outcome, the predictor's placement, and the agent's resulting action. The requirement that the predictor be near-certainly correct is the requirement that, across the consistent configurations of this joint system, prediction and action agree. Configurations in which they diverge---the coin sends the agent to two boxes while the predictor placed the million, or vice versa---are exactly the configurations the agent would, on opening the boxes, experience as the predictor having been wrong; and a near-perfect predictor is, by hypothesis, almost never wrong. Those divergent configurations are therefore of small positive weight, and what remains is a \emph{coupling} between the coin outcome and the placement: not a causal influence of either on the other, and not a signal, but a constraint on which joint configurations are consistent with the stipulated near-certainty. The random device does not break the correlation; it is absorbed into it.

\textit{The same mechanism, and the retrocausal air dissolved.} This is the same joint-consistency mechanism that the broader framework uses to handle correlations across events with no classical common cause \cite{Aethus}, applied here to a decision-theoretic rather than a physical setting; the framework's own detailed development treats this supercomputer-and-coin configuration as its illustrative case \cite{Aethus}. The apparent retrocausal air of the result---how can the coin ``know'' what was placed?---dissolves once one sees that nothing is transmitted and nothing is caused: with a perfect predictor the mismatch configurations are jointly invalid and absent, while with positive error they are possible and carry small weight; invalidity is reserved for the exact exclusion.



Quantitatively, the mechanism and the resolve-gradient of Section~\ref{sec:proposal} are the same formula. A deliberate randomizer at bias $r$, under this appendix's stated extra assumption that the accuracy entry holds against randomizers, has exactly the $\mathrm{EU}(r)$ displayed with the decision rule --- continuous in $r$ and matching the pure verdicts at the endpoints. Nozick's stipulation that randomizers find the opaque box empty instead gives $\mathrm{EU}_{N}(r) = (1-r)\,t$, which jumps discontinuously at both endpoints --- by $(1-\varepsilon)M$ at $r = 1$ and by $\varepsilon M$ at $r = 0$. Different treatment of a randomized and a deterministic policy is not automatically defective, since the predictor may respond to the policy category; continuity is an attractive modeling property under the paper's assumptions, not an independent justification for them. The framework's treatment is thus not merely mechanism rather than stipulation: it is a continuous completion of the two pure cases, where the stipulation is discontinuous at both.

\printbibliography
\end{document}